% Generated by the blueprint tool (blueprint-pdf/1) from a blueprint snapshot; edit the
% Markdown sources, not this file.  Compile with `latexmk -pdf` or by running
% `pdflatex` twice (the second run settles the contents and the page references).
% Images, if any, are read from assets/ next to this file.
\documentclass[11pt]{article}
\usepackage[T1]{fontenc}
\usepackage[utf8]{inputenc}
\usepackage{lmodern}
\usepackage{textcomp}
\usepackage[a4paper,margin=2.5cm,headheight=14pt]{geometry}
\usepackage{amsmath,amssymb,amsthm}
\usepackage{amscd}  % KaTeX's CD diagrams; part of amsmath, so always there
\usepackage{mathtools}
\usepackage{bm}
\usepackage{graphicx}
\usepackage{xcolor}
\usepackage{array,booktabs,tabularx}
\usepackage{fancyvrb}
\usepackage[normalem]{ulem}
\usepackage{microtype}
% Optional, and only with their fonts: a small TeX installation has the
% packages without the fonts, which fails at the first formula.
\IfFileExists{mathrsfs.sty}{\IfFileExists{rsfs10.tfm}{\usepackage{mathrsfs}}{}}{}
\IfFileExists{stmaryrd.sty}{\IfFileExists{stmary10.tfm}{\usepackage{stmaryrd}}{}}{}
\IfFileExists{cancel.sty}{\usepackage{cancel}}{}
\usepackage[hyphens]{url}
\usepackage[pdfusetitle,bookmarksnumbered=false]{hyperref}
\usepackage{bookmark}
\definecolor{bplink}{HTML}{1F4E9C}
\definecolor{bpmuted}{HTML}{5F676E}
\hypersetup{colorlinks=true,linkcolor=bplink,urlcolor=bplink,citecolor=bplink}
\setlength{\parindent}{0pt}
\setlength{\parskip}{0.55em plus 0.2em minus 0.1em}
% Lists as tight as the page's.
\makeatletter
\def\@listi{\leftmargin\leftmargini\parsep 0pt\topsep 0.3em plus 0.1em\itemsep 0.1em plus 0.1em}
\let\@listI\@listi
\def\@listii{\leftmargin\leftmarginii\labelwidth\leftmarginii\advance\labelwidth-\labelsep\parsep 0pt\topsep 0.1em\itemsep 0.1em}
% An ordered list that does not start at 1.
\newcommand{\bpstartat}[1]{\setcounter{\@enumctr}{\numexpr#1-1\relax}}
\makeatother
\setcounter{tocdepth}{3}
\emergencystretch=3em
\allowdisplaybreaks

% A numbered statement, set as a paper sets one: "Theorem 2.1 (Title)."
% #1 is the lead ("Theorem~2.1"), #2 what follows it in roman.
\newenvironment{bpstatement}[2]{\par\addvspace{0.6\baselineskip}\noindent\textbf{#1}#2\textbf{.}\enspace\ignorespaces}
  {\par\addvspace{0.4\baselineskip}}
\newenvironment{bpstatementit}[2]{\begin{bpstatement}{#1}{#2}\itshape\ignorespaces}{\end{bpstatement}}
% A step: a relation with prose of its own, read under the object it starts from.
\newenvironment{bpstep}[1]{\par\addvspace{0.3\baselineskip}\begin{list}{}{\leftmargin=1.6em\rightmargin=0pt\topsep=0pt}\item\relax
  {\small\color{bpmuted}#1}\par\nopagebreak}{\end{list}\par}
% Notes about where an object is written out.
\newenvironment{bpnote}{\par\small\color{bpmuted}\itshape}{\par}
\newcommand{\bpomitted}[1]{\par{\small\color{bpmuted}\itshape #1\par}}
\newcommand{\bpbodyheading}[1]{\par\addvspace{0.4\baselineskip}\noindent\textbf{#1}\par\nopagebreak}
\newcommand{\bprule}{\par\medskip\noindent\hfil\rule{0.3\linewidth}{0.4pt}\hfil\par\medskip}
\newcommand{\bpboundary}[1]{\par\nopagebreak{\small\color{bpmuted}#1\par}\nopagebreak}
% A reference to something this export leaves out, and one that resolves nowhere.
\newcommand{\bpout}[1]{#1\textsuperscript{\hyperlink{bp-outside}{\ensuremath{\ast}}}}
\newcommand{\bpbroken}[1]{\textcolor{red!70!black}{[#1?]}}
\newcommand{\bpslug}[1]{\texttt{#1}}
% A bibliography entry: label, anchor, entry; the label hangs in the margin.
\newcommand{\bpbibitem}[3]{\par\setbox0\hbox{[#1]\enspace}\ifdim\wd0<5em\setbox0\hbox to 5em{[#1]\hfil}\fi
  \hangindent=\wd0 \hangafter=1 \noindent\box0\phantomsection\label{#2}#3\par\smallskip}
% A character the export knows no way to set.
\newcommand{\bpglyph}[1]{\fbox{\scriptsize #1}}
% Images, at their own size unless they are wider than the line.
\newsavebox{\bpimagebox}
\newcommand{\bpimage}[1]{\sbox{\bpimagebox}{\includegraphics{#1}}%
  \ifdim\wd\bpimagebox>\linewidth\resizebox{\linewidth}{!}{\usebox{\bpimagebox}}\else\usebox{\bpimagebox}\fi}
\newcommand{\bpmissingimage}[1]{\fbox{\parbox{0.8\linewidth}{\centering\small\color{bpmuted}Image not included: #1}}}
\DefineVerbatimEnvironment{bpverbatim}{Verbatim}{fontsize=\small,xleftmargin=1em}
\renewcommand{\qedsymbol}{\ensuremath{\blacksquare}}
\hypersetup{pdftitle={Weil Conjectures — 9 Smooth purity},pdfsubject={Chapter 9 Smooth purity; all levels},pdfcreator={blueprint}}
% KaTeX commands the website knows
\providecommand{\R}{\mathbb{R}}\providecommand{\N}{\mathbb{N}}\providecommand{\Z}{\mathbb{Z}}
\providecommand{\Q}{\mathbb{Q}}\providecommand{\C}{\mathbb{C}}
\providecommand{\Reals}{\mathbb{R}}\providecommand{\reals}{\mathbb{R}}\providecommand{\Complex}{\mathbb{C}}
\providecommand{\cnums}{\mathbb{C}}\providecommand{\natnums}{\mathbb{N}}\providecommand{\Bbb}[1]{\mathbb{#1}}
\providecommand{\bold}[1]{\mathbf{#1}}\providecommand{\frak}[1]{\mathfrak{#1}}\providecommand{\mathbbm}[1]{\mathbb{#1}}
\providecommand{\lang}{\langle}\providecommand{\rang}{\rangle}\providecommand{\bra}[1]{\langle #1\rvert}
\providecommand{\ket}[1]{\lvert #1\rangle}\providecommand{\braket}[1]{\langle #1\rangle}
\providecommand{\set}[1]{\{#1\}}\providecommand{\Set}[1]{\left\{#1\right\}}
\providecommand{\infin}{\infty}\providecommand{\isin}{\in}\providecommand{\plusmn}{\pm}\providecommand{\sub}{\subset}
\providecommand{\sube}{\subseteq}\providecommand{\supe}{\supseteq}\providecommand{\exist}{\exists}
\providecommand{\bull}{\bullet}\providecommand{\larr}{\leftarrow}\providecommand{\rarr}{\rightarrow}
\providecommand{\lrarr}{\leftrightarrow}\providecommand{\Larr}{\Leftarrow}\providecommand{\Rarr}{\Rightarrow}
\providecommand{\Lrarr}{\Leftrightarrow}\providecommand{\uarr}{\uparrow}\providecommand{\darr}{\downarrow}
\providecommand{\harr}{\leftrightarrow}\providecommand{\lArr}{\Leftarrow}\providecommand{\rArr}{\Rightarrow}
\providecommand{\hArr}{\Leftrightarrow}\providecommand{\alef}{\aleph}\providecommand{\alefsym}{\aleph}
\providecommand{\real}{\Re}\providecommand{\image}{\Im}\providecommand{\weierp}{\wp}
\providecommand{\Alpha}{\mathrm{A}}\providecommand{\Beta}{\mathrm{B}}\providecommand{\Epsilon}{\mathrm{E}}
\providecommand{\Zeta}{\mathrm{Z}}\providecommand{\Eta}{\mathrm{H}}\providecommand{\Iota}{\mathrm{I}}
\providecommand{\Kappa}{\mathrm{K}}\providecommand{\Mu}{\mathrm{M}}\providecommand{\Nu}{\mathrm{N}}
\providecommand{\Omicron}{\mathrm{O}}\providecommand{\Rho}{\mathrm{P}}\providecommand{\Tau}{\mathrm{T}}
\providecommand{\Chi}{\mathrm{X}}\providecommand{\omicron}{o}\providecommand{\gt}{>}\providecommand{\lt}{<}
\providecommand{\And}{\&}\providecommand{\hdashline}{\hline}\providecommand{\operatornamewithlimits}[1]{\operatorname*{#1}}
\providecommand{\llbracket}{[\![}\providecommand{\rrbracket}{]\!]}\providecommand{\cancel}[1]{#1}
\providecommand{\mathscr}{\mathcal}
% the characters beyond ASCII this document uses
\DeclareUnicodeCharacter{2160}{\textup{I}}
\DeclareUnicodeCharacter{2161}{\textup{II}}
% the project's macros ([katex.macros] in blueprint.toml)
\providecommand{\A}{}\renewcommand{\A}{\mathbf A}
\providecommand{\AEn}{}\renewcommand{\AEn}{A_{E,n}}
\providecommand{\Aaf}{}\renewcommand{\Aaf}{\mathbf A}
\providecommand{\Ab}{}\renewcommand{\Ab}{\mathsf{Ab}}
\providecommand{\AffSch}{}\renewcommand{\AffSch}{\operatorname{AffSch}}
\providecommand{\Aone}{}\renewcommand{\Aone}{\mathbf A^1}
\providecommand{\Bl}{}\renewcommand{\Bl}{\operatorname{Bl}}
\providecommand{\Br}{}\renewcommand{\Br}{\mathrm{Br}}
\providecommand{\C}{}\renewcommand{\C}{\mathbb C}
\providecommand{\Cech}{}\renewcommand{\Cech}{\check{H}}
\providecommand{\Coext}{}\renewcommand{\Coext}{\operatorname{Coext}}
\providecommand{\Cone}{}\renewcommand{\Cone}{\operatorname{Cone}}
\providecommand{\Cons}{}\renewcommand{\Cons}{\operatorname{Cons}}
\providecommand{\Cref}{}\renewcommand{\Cref}[1]{\text{#1}}
\providecommand{\D}{}\renewcommand{\D}{D}
\providecommand{\DLAFT}{}\renewcommand{\DLAFT}{\operatorname{FT}}
\providecommand{\DLAFr}{}\renewcommand{\DLAFr}{\operatorname{Fr}}
\providecommand{\DLAQ}{}\renewcommand{\DLAQ}{{\overline{\mathbf Q}_\ell}}
\providecommand{\DLRQ}{}\renewcommand{\DLRQ}{{\overline{\mathbf Q}_\ell}}
\providecommand{\Dbc}{}\renewcommand{\Dbc}{D^b_c}
\providecommand{\Dc}{}\renewcommand{\Dc}{D_{\mathrm{cons}}}
\providecommand{\Dctf}{}\renewcommand{\Dctf}{D^b_{ctf}}
\providecommand{\Dic}{}\renewcommand{\Dic}{D_{\IndCons}}
\providecommand{\Dplus}{}\renewcommand{\Dplus}{D^{+}}
\providecommand{\End}{}\renewcommand{\End}{\operatorname{End}}
\providecommand{\Et}{}\renewcommand{\Et}{\mathrm{\acute{e}t}}
\providecommand{\Ex}{}\renewcommand{\Ex}{\mathrm{Ex}}
\providecommand{\Ext}{}\renewcommand{\Ext}{\operatorname{Ext}}
\providecommand{\F}{}\renewcommand{\F}{\mathbb F}
\providecommand{\FT}{}\renewcommand{\FT}{\operatorname{FT}}
\providecommand{\Fbar}{}\renewcommand{\Fbar}{{\overline{\mathbf F}_q}}
\providecommand{\Fell}{}\renewcommand{\Fell}{\mathbf F_\ell}
\providecommand{\Fib}{}\renewcommand{\Fib}{\operatorname{Fib}}
\providecommand{\Fix}{}\renewcommand{\Fix}{\operatorname{Fix}}
\providecommand{\Fp}{}\renewcommand{\Fp}{{\mathbf F_{p}}}
\providecommand{\Fq}{}\renewcommand{\Fq}{{\mathbf F_q}}
\providecommand{\Fqn}{}\renewcommand{\Fqn}{{\mathbf F_{q^{n}}}}
\providecommand{\Frob}{}\renewcommand{\Frob}{\operatorname{Frob}}
\providecommand{\GL}{}\renewcommand{\GL}{\operatorname{GL}}
\providecommand{\Gal}{}\renewcommand{\Gal}{\operatorname{Gal}}
\providecommand{\Hens}{}\renewcommand{\Hens}{\operatorname{Hens}}
\providecommand{\Hilb}{}\renewcommand{\Hilb}{\operatorname{Hilb}}
\providecommand{\Hom}{}\renewcommand{\Hom}{\operatorname{Hom}}
\providecommand{\IndCons}{}\renewcommand{\IndCons}{\operatorname{IndCons}}
\providecommand{\Kar}{}\renewcommand{\Kar}{\operatorname{Kar}}
\providecommand{\LT}{}\renewcommand{\LT}{\operatorname{LT}}
\providecommand{\Lie}{}\renewcommand{\Lie}{\operatorname{Lie}}
\providecommand{\Loc}{}\renewcommand{\Loc}{\operatorname{Loc}}
\providecommand{\Mod}{}\renewcommand{\Mod}{\mathsf{Mod}}
\providecommand{\N}{}\renewcommand{\N}{\mathbb{N}}
\providecommand{\Needspace}{}\renewcommand{\Needspace}{\par}
\providecommand{\O}{}\renewcommand{\O}{{\mathcal{O}}}
\providecommand{\Ob}{}\renewcommand{\Ob}{\mathrm{Ob}}
\providecommand{\Oe}{}\renewcommand{\Oe}{\mathcal O_E}
\providecommand{\PP}{}\renewcommand{\PP}{\mathbf{P}}
\providecommand{\Perf}{}\renewcommand{\Perf}{\operatorname{Perf}}
\providecommand{\Pic}{}\renewcommand{\Pic}{\mathrm{Pic}}
\providecommand{\Pm}{}\renewcommand{\Pm}{\mathbf P^{m}}
\providecommand{\Pone}{}\renewcommand{\Pone}{\mathbf P^{1}}
\providecommand{\Pqn}{}\renewcommand{\Pqn}{\mathbf P^{n}}
\providecommand{\ProFin}{}\renewcommand{\ProFin}{\operatorname{ProFin}}
\providecommand{\Proet}{}\renewcommand{\Proet}{\mathrm{pro\acute{e}t}}
\providecommand{\Q}{}\renewcommand{\Q}{\mathbb Q}
\providecommand{\Qb}{}\renewcommand{\Qb}{\overline{\mathbf Q}_{\ell}}
\providecommand{\Qbar}{}\renewcommand{\Qbar}{\overline{\mathbf Q}}
\providecommand{\QbarX}{}\renewcommand{\QbarX}{\overline{\mathbf Q}_{\ell,X}}
\providecommand{\Qbarl}{}\renewcommand{\Qbarl}{{\overline{\mathbf Q}_{\ell}}}
\providecommand{\Qell}{}\renewcommand{\Qell}{{\mathbf{Q}_{\ell}}}
\providecommand{\Ql}{}\renewcommand{\Ql}{\mathbf Q_{\ell}}
\providecommand{\Qlb}{}\renewcommand{\Qlb}{{\overline{\mathbf Q}_\ell}}
\providecommand{\Qlbar}{}\renewcommand{\Qlbar}{\overline{\mathbf Q}_\ell}
\providecommand{\R}{}\renewcommand{\R}{\mathbb{R}}
\providecommand{\RG}{}\renewcommand{\RG}{\operatorname{\mathrm{R}\Gamma}}
\providecommand{\RGamma}{}\renewcommand{\RGamma}{R\Gamma}
\providecommand{\RGammac}{}\renewcommand{\RGammac}{R\Gamma_c}
\providecommand{\RHom}{}\renewcommand{\RHom}{R\mathcal{H}om}
\providecommand{\RHomX}{}\renewcommand{\RHomX}{R\mathcal H om_X}
\providecommand{\Rat}{}\renewcommand{\Rat}{\operatorname{IsRat}}
\providecommand{\Res}{}\renewcommand{\Res}{\operatorname{Res}}
\providecommand{\Rlim}{}\renewcommand{\Rlim}{R\!\varprojlim}
\providecommand{\SIm}{}\renewcommand{\SIm}{\operatorname{SIm}}
\providecommand{\Sch}{}\renewcommand{\Sch}{\mathsf{Sch}}
\providecommand{\Sh}{}\renewcommand{\Sh}{\operatorname{Sh}}
\providecommand{\Span}{}\renewcommand{\Span}{\operatorname{span}}
\providecommand{\Spec}{}\renewcommand{\Spec}{\operatorname{Spec}}
\providecommand{\Tr}{}\renewcommand{\Tr}{\operatorname{Tr}}
\providecommand{\Trsm}{}\renewcommand{\Trsm}{\operatorname{Tr}^{\mathrm{sm}}}
\providecommand{\Wt}{}\renewcommand{\Wt}{\operatorname{Wt}}
\providecommand{\Z}{}\renewcommand{\Z}{\mathbb{Z}}
\providecommand{\Zar}{}\renewcommand{\Zar}{\mathrm{zar}}
\providecommand{\Zell}{}\renewcommand{\Zell}{{\mathbf{Z}_{\ell}}}
\providecommand{\Zl}{}\renewcommand{\Zl}{{\mathbf Z_\ell}}
\providecommand{\affet}{}\renewcommand{\affet}[1]{{#1}^{\mathrm{aff}}_{\mathrm{ét}}}
\providecommand{\affproet}{}\renewcommand{\affproet}[1]{{#1}^{\mathrm{aff}}_{\mathrm{proét}}}
\providecommand{\boxt}{}\renewcommand{\boxt}{\boxtimes_A^{\mathrm L}}
\providecommand{\cF}{}\renewcommand{\cF}{\mathcal{F}}
\providecommand{\cHom}{}\renewcommand{\cHom}{R\mathcal{H}\!om}
\providecommand{\cI}{}\renewcommand{\cI}{\mathcal{I}}
\providecommand{\cL}{}\renewcommand{\cL}{\mathcal L}
\providecommand{\cM}{}\renewcommand{\cM}{\mathcal M}
\providecommand{\cN}{}\renewcommand{\cN}{\mathcal N}
\providecommand{\cO}{}\renewcommand{\cO}{\mathcal O}
\providecommand{\cT}{}\renewcommand{\cT}{\mathcal T}
\providecommand{\checked}{}\renewcommand{\checked}{\color{red}\aftergroup\restorecolor}
\providecommand{\coeff}{}\renewcommand{\coeff}{\operatorname{coeff}}
\providecommand{\coker}{}\renewcommand{\coker}{\operatorname{coker}}
\providecommand{\colim}{}\renewcommand{\colim}{\operatorname{colim}}
\providecommand{\constructedby}{}\renewcommand{\constructedby}[1]{\par\smallskip\noindent\textcolor{purple!60!black}{\textsf{Constructed by:} }
  \footnotesize\sloppy\nolinkurl{#1}\par}
\providecommand{\deepresult}{}\renewcommand{\deepresult}[1]{\par\smallskip\noindent\textcolor{orange!45!black}{
  \textsf{Blueprint role: named deep theorem node; see the cluster proof architecture.}}
  \par\noindent\textcolor{orange!45!black}{\textsf{Mathematical source:} }#1\par}
\providecommand{\dt}{}\renewcommand{\dt}{\mathbin{\otimes^{\mathbf L}}}
\providecommand{\dw}{}\renewcommand{\dw}{\operatorname{dw}}
\providecommand{\et}{}\renewcommand{\et}[1]{{#1}_{\mathrm{ét}}}
\providecommand{\ev}{}\renewcommand{\ev}{\operatorname{ev}}
\providecommand{\fib}{}\renewcommand{\fib}{\operatorname{fib}}
\providecommand{\formalizationnote}{}\renewcommand{\formalizationnote}[1]{\begin{quote}\small\raggedright\textbf{Formalization note.} #1\end{quote}}
\providecommand{\fpqc}{}\renewcommand{\fpqc}{X_{\mathrm{fpqc}}}
\providecommand{\gammacont}{}\renewcommand{\gammacont}{\Gamma_{\mathrm{cont}}}
\providecommand{\gm}{}\renewcommand{\gm}{\mathbf G_m}
\providecommand{\hocolim}{}\renewcommand{\hocolim}{\operatorname{hocolim}}
\providecommand{\holim}{}\renewcommand{\holim}{\operatorname{holim}}
\providecommand{\hrB}{}\renewcommand{\hrB}{\@ifnextchar]{\hrEnd}{\hrStep}}
\providecommand{\hrEnd}{}\renewcommand{\hrEnd}{\@firstoftwo{}}
\providecommand{\hrStep}{}\renewcommand{\hrStep}{\@firstoftwo{\hrB}}
\providecommand{\id}{}\renewcommand{\id}{\operatorname{id}}
\providecommand{\kunFl}{}\renewcommand{\kunFl}{\mathbf F_\ell}
\providecommand{\kunZn}{}\renewcommand{\kunZn}{\mathbf Z/n\mathbf Z}
\providecommand{\kuncd}{}\renewcommand{\kuncd}{\operatorname{cd}}
\providecommand{\kundim}{}\renewcommand{\kundim}{\dim}
\providecommand{\kunotimes}{}\renewcommand{\kunotimes}[1]{\otimes^{\mathbf L}_{#1}}
\providecommand{\kunsupp}{}\renewcommand{\kunsupp}{\operatorname{Supp}}
\providecommand{\lotimes}{}\renewcommand{\lotimes}{\otimes^{\mathrm L}}
\providecommand{\m}{}\renewcommand{\m}{{\mathfrak{m}}}
\providecommand{\n}{}\renewcommand{\n}{{\mathfrak{n}}}
\providecommand{\nagataGap}{}\renewcommand{\nagataGap}[1]{\par\noindent\textbf{Formalization gap.} #1\par}
\providecommand{\openobligation}{}\renewcommand{\openobligation}[1]{\par\smallskip\noindent\textcolor{red!65!black}{\textsf{Open obligation:} }
  \footnotesize\sloppy\nolinkurl{#1}\par}
\providecommand{\p}{}\renewcommand{\p}{{\mathfrak{p}}}
\providecommand{\pr}{}\renewcommand{\pr}{\operatorname{pr}}
\providecommand{\preshv}{}\renewcommand{\preshv}[1]{\mathrm{PreShv}(#1)}
\providecommand{\pro}{}\renewcommand{\pro}[1]{\mathrm{Pro}(#1)}
\providecommand{\proet}{}\renewcommand{\proet}{X_{\mathrm{proét}}}
\providecommand{\q}{}\renewcommand{\q}{{\mathfrak{q}}}
\providecommand{\rad}{}\renewcommand{\rad}{\operatorname{rad}}
\providecommand{\rank}{}\renewcommand{\rank}{\operatorname{rank}}
\providecommand{\red}{}\renewcommand{\red}{\color{red}}
\providecommand{\restorecolor}{}\renewcommand{\restorecolor}{\color{black}}
\providecommand{\retiredobligation}{}\renewcommand{\retiredobligation}[1]{\par\smallskip\noindent\textcolor{red!75!black}{\textsf{Retired invalid route:} }
  \footnotesize\sloppy\nolinkurl{#1}\par}
\providecommand{\rk}{}\renewcommand{\rk}{\operatorname{rk}}
\providecommand{\sep}{}\renewcommand{\sep}{\textnormal{sep}}
\providecommand{\shv}{}\renewcommand{\shv}[1]{\mathrm{Shv}(#1)}
\providecommand{\spec}{}\renewcommand{\spec}[1]{\mathrm{Spec}(#1)}
\providecommand{\stacks}{}\renewcommand{\stacks}[1]{\href{https://stacks.math.columbia.edu/tag/#1}{\textnormal{Stacks, Tag~#1}}}
\providecommand{\stacksproject}{}\renewcommand{\stacksproject}[1]{{\cite[\href{https://stacks.math.columbia.edu/tag/#1}{Tag #1}]{stacks-project}}}
\providecommand{\statusnote}{}\renewcommand{\statusnote}[1]{\begin{quote}\small\raggedright\textbf{Status.} #1\end{quote}}
\providecommand{\tildeideal}{}\renewcommand{\tildeideal}[2]{\widetilde{(#1, #2)}}
\providecommand{\tildering}{}\renewcommand{\tildering}[3]{{#1}_{\widetilde{(#2, #3)}}}
\providecommand{\tr}{}\renewcommand{\tr}{\operatorname{tr}}
\providecommand{\wt}{}\renewcommand{\wt}{\operatorname{wt}}

\begin{document}
\pagestyle{plain}
\begin{center}
{\LARGE\bfseries Weil Conjectures\par}
\vspace{0.7em}
{\large Chapter 9: Smooth purity\par}
\vspace{0.5em}
{\small\color{bpmuted}Snapshot 2026-09-26, revision 8e4cb88e2a\par}
\end{center}
\vspace{0.5em}
{\small\tableofcontents}
\bigskip\hrule\bigskip

% object: chap-Smooth-purity (blueprint/chap-Smooth-purity/_section.md)
\section*{\texorpdfstring{9\quad Smooth purity}{9 Smooth purity}}
\phantomsection\label{bp1}%
\addcontentsline{toc}{section}{\texorpdfstring{9\quad Smooth purity}{9 Smooth purity}}
% text: prose-09-01
\phantomsection\label{bp2}%
Poincaré duality for a smooth variety \(X\) of dimension \(d\) pairs \(H^i_c(X,\Lambda)\) with \(H^{2d-i}(X,\Lambda(d))\) through the cup product and a trace map \(H^{2d}_c(X,\Lambda(d))\to\Lambda\). If \(a\) is the structure map of \(X\), such a trace map is the counit of \(Ra_!\dashv a^!\), once we know that \(a^!\Lambda\simeq\Lambda_X(d)[2d]\). But Part I produces \(f^!\) (Theorem \bpout{5.2}) as a right adjoint by Brown representability, with no formula.

A formula exists on charts. Consider \(f:X\to Y\) smooth of pure relative dimension \(d\), and an open subscheme \(U\subset X\) on which \(f\) factors as an étale map \(e:U\to\mathbf A^d_Y\) followed by the projection \(p_d\). With finite coefficients \(f^!K|_U\simeq e^*p_d^*K(d)[2d]\) (Proposition \bpout{6.9}), and under this isomorphism the counit is integration along \(d\) affine lines. This integration comes from the trace map of \(\mathbf P^1\) (Proposition \bpout{6.7}), which sends the class of a point to \(1\). The question is then whether these chart descriptions agree on overlaps.

In topology they need not. For a submersion \(f:X\to Y\) of manifolds with fibres of real dimension \(m\), Verdier's theory gives \(f^!K\simeq f^*K\otimes\mathrm{or}_{X/Y}[m]\). Here the relative orientation sheaf \(\mathrm{or}_{X/Y}\) is locally constant of rank one, and it is trivial exactly when \(f\) can be oriented. A linear change of coordinates \(g\) acts on \(H^m_c(\mathbf R^m)\) by the sign of \(\det g\), and since \(\operatorname{GL}_m(\mathbf R)\) has two components, the local identifications need not glue. For complex manifolds, by contrast, \(\operatorname{GL}_d(\mathbf C)\) is connected and the orientation is canonical.

The same argument works for us on a chart. Indeed, a linear change of coordinates multiplies the trace map of \(\mathbf A^d\) by a locally constant unit on \(\operatorname{GL}_d\), and this unit is \(1\) at the identity, hence everywhere, since \(\operatorname{GL}_d\) has connected geometric fibres. So every vector bundle has a canonical trace map (Proposition \hyperref[bp19]{9.6}). The trace map has source \(H^{2d}_c(-,\Lambda(d))\), so the Tate twist \(\Lambda(d)\) plays the role of the orientation sheaf. Over an algebraically closed field a compatible system of roots of unity trivializes \(\Lambda(d)\), but not canonically, and over \(\mathbf F_q\) the twist is visible: there geometric Frobenius acts on \(\Lambda(1)\) by \(q^{-1}\) (Section \bpout{10.5}), and this is where the factor \(q^d\) in Weil's Poincaré duality comes from.

Connectedness of \(\operatorname{GL}_d\) does not settle the gluing, because two étale charts \(U\to\mathbf A^d_Y\) differ by a map that need not be linear. We therefore construct the isomorphism without charts, from the diagonal, in the way a tubular neighbourhood linearizes a submanifold (Section \hyperref[bp22]{9.4}).

The outcome, Theorem \hyperref[bp9]{9.3}, is a canonical isomorphism

\[
\gamma_f:f^*K(d)[2d]\xrightarrow{\sim}f^!K,
\]

natural in \(K\in D^b_{\mathrm{ctf}}(Y,\Lambda)\) for every admissible \(\Lambda\). It is compatible with base change, composition and change of coefficients, and it is normalized by the trace map of the affine line. This statement is Deligne's [\hyperref[bp99]{SGA4}, XVIII], where it is called Poincaré duality.

% object: sec-Statement (blueprint/chap-Smooth-purity/sec-Statement/_section.md), content hidden
\subsection*{\texorpdfstring{9.1\quad Statement}{9.1 Statement}}
\phantomsection\label{bp3}%
\addcontentsline{toc}{subsection}{\texorpdfstring{9.1\quad Statement}{9.1 Statement}}
% text: prose-09-02
\phantomsection\label{bp4}%
Smooth purity is stated for six kinds of coefficients at once, because the proof moves between them: the geometry is done with finite rings, the case \(\mathcal O_E\) is a limit of finite levels, and the fields are obtained by inverting \(\ell\). Throughout this chapter \(\Lambda\) is an admissible coefficient ring, that is, a finite commutative ring killed by a power of \(\ell\), \(\mathcal O_E\), \(E\) or \(\overline{\mathbf Q}_\ell\) (the rings \(\mathbf Z_\ell\) and \(\mathbf Q_\ell\) being the cases \(E=\mathbf Q_\ell\)), and Tate twists are those of Definition \bpout{2.53}.

The theorem asserts one isomorphism on all of \(X\), whereas on a chart an isomorphism of this shape is already given by Proposition \bpout{6.9} with finite coefficients. A bare isomorphism for each \(f\) separately, however, would not be enough for what follows, and the theorem therefore also lists three properties that the later chapters use. Property (2), compatibility with base change, restricts \(\gamma_f\) to the fibres of \(f\), and it is what makes \(\gamma_f\) commute with Frobenius, because the Frobenius square is a base change (Proposition \bpout{10.20}). Property (3), compatibility with composition, gives the transitivity of the trace map (Proposition \hyperref[bp86]{9.26}), the analogue of integrating over the fibres of a composite in two steps. Property (4), compatibility with change of coefficients, lets us construct \(\gamma_f\) with finite coefficients and pass to \(\mathcal O_E\), \(E\) and \(\overline{\mathbf Q}_\ell\), and it also makes Poincaré duality compatible with change of coefficients (Theorem \bpout{10.12}). Properties (1) and (2) involve maps that we define now, for every admissible \(\Lambda\); the functors are those of Theorem \bpout{5.2}, restricted to \(D^b_{\mathrm{ctf}}\).

The first map compares \(f^!\) before and after a change of base. Let \(f:X\to Y\) and \(g:Y'\to Y\) be morphisms of separated schemes of finite type over \(k\), let \(X'=X\times_YY'\), and let \(f':X'\to Y'\) and \(g':X'\to X\) be the projections. Write \(\varepsilon_f(K):Rf_!f^!K\to K\) for the counit of \(Rf_!\dashv f^!\).

% object: definition-exceptional-base-change-map (blueprint/chap-Smooth-purity/sec-Statement/definition-exceptional-base-change-map/_section.md)
\phantomsection\label{bp5}%
\begin{bpstatement}{Definition~9.1}{ (Exceptional base change map)}
For \(K\in D^b_{\mathrm{ctf}}(Y,\Lambda)\), the \emph{exceptional base change map} \(g'^*f^!K\to f'^!g^*K\) is the adjoint transpose, for \(Rf'_!\dashv f'^!\), of the composite

\[
Rf'_!g'^*f^!K\xrightarrow{\sim}g^*Rf_!f^!K
\xrightarrow{\ g^*\varepsilon_f(K)\ }g^*K,
\]

whose first arrow is the inverse of the base change isomorphism of Proposition proposition-composition-base-change-rf-shriek.
\end{bpstatement}

% text: prose-09-03
\phantomsection\label{bp6}%
The second map is integration along the affine line. Let \(p:\mathbf A^1_Y\to Y\) be the projection, and factor it as \(p=q\circ j\), where \(j:\mathbf A^1_Y\hookrightarrow\mathbf P^1_Y\) is the complement of the section at infinity and \(q:\mathbf P^1_Y\to Y\) is the projection.

% object: definition-trace-map-affine-line (blueprint/chap-Smooth-purity/sec-Statement/definition-trace-map-affine-line/_section.md)
\phantomsection\label{bp7}%
\begin{bpstatement}{Definition~9.2}{ (Trace map of the affine line)}
The \emph{trace map of the affine line} is the composite

\[
\operatorname{Tr}_p:Rp_!\Lambda_{\mathbf A^1_Y}(1)[2]
\simeq Rq_*j_!\Lambda_{\mathbf A^1_Y}(1)[2]
\longrightarrow Rq_*\Lambda_{\mathbf P^1_Y}(1)[2]
\xrightarrow{\ \operatorname{Tr}_q\ }\Lambda_Y,
\tag{9.1}
\]

where the first arrow is induced by the counit \(j_!j^*\to\mathrm{id}\) and \(\operatorname{Tr}_q\) is the trace map of the projective line: for finite \(\Lambda\) it is the map (6.5) of Lemma \bpout{6.4} for the projection \(q\), and otherwise it is the trace of Proposition \bpout{6.22}.
\end{bpstatement}

% text: prose-09-04
\phantomsection\label{bp8}%
Over an algebraically closed field the trace map of the affine line sends the class in \(H^2_c(\mathbf A^1,\Lambda(1))\) of a point to \(1\), as we check with finite coefficients in Section \hyperref[bp15]{9.3}.

% object: theorem-smooth-purity-all-admissible-coefficients (blueprint/chap-Smooth-purity/sec-Statement/theorem-smooth-purity-all-admissible-coefficients/_section.md)
\phantomsection\label{bp9}%
\begin{bpstatementit}{Theorem~9.3}{ (Smooth purity for all admissible coefficients)}
Let \(f:X\to Y\) be smooth of pure relative dimension \(d\). For every admissible coefficient ring \(\Lambda\in\{A,\mathbf Z_\ell,\mathcal O_E,\mathbf Q_\ell,E,\overline{\mathbf Q}_\ell\}\), with \(A\) finite, there are isomorphisms, natural in \(K\in D^b_{\mathrm{ctf}}(Y,\Lambda)\),

\[
\gamma_f(K):
f^{*}K(d)[2d]\xrightarrow{\sim}f^{!}K,
\tag{9.2}
\]

with the following properties.

\begin{enumerate}
\item If \(f\) is the projection \(p:\mathbf A^1_Y\to Y\), then \(\varepsilon_p(\Lambda_Y)\circ Rp_!\bigl(\gamma_p(\Lambda_Y)\bigr)\) is the trace map (9.1).
\item For \(g:Y'\to Y\) as in Definition \hyperref[bp5]{9.1}, with exceptional base change map \(\beta:g'^*f^!K\to f'^!g^*K\), we have \(\beta\circ g'^*\gamma_f(K)=\gamma_{f'}(g^*K)\), the sources being identified by \(g'^*f^*\simeq f'^*g^*\).
\item For a smooth morphism \(h:Y\to Z\) of pure relative dimension \(e\) and \(K\in D^b_{\mathrm{ctf}}(Z,\Lambda)\), the map \(\gamma_{h\circ f}(K)\) is the composite of \(f^*(\gamma_h(K))(d)[2d]\) and \(\gamma_f(h^!K)\), the sources being identified by \((h\circ f)^*\simeq f^*h^*\) and the targets by the composition isomorphism \(f^!h^!\simeq(h\circ f)^!\) of Proposition proposition-composition-open-restriction-compactifications-f.
\item For a coefficient map \(\Lambda\to\Lambda'\) of Proposition \bpout{7.10}, the change-of-coefficients isomorphism \((f^!K)_{\Lambda'}\to f^!(K_{\Lambda'})\) of that proposition carries \((\gamma_{f,\Lambda}(K))_{\Lambda'}\) to \(\gamma_{f,\Lambda'}(K_{\Lambda'})\), the sources being identified by the change-of-coefficients isomorphism for \(f^*\).
\end{enumerate}
\end{bpstatementit}

% text: prose-09-05
\phantomsection\label{bp10}%
We recall what \(D^b_{\mathrm{ctf}}\) means for each coefficient ring. For integral coefficients and for coefficients in a field, \(D^b_{\mathrm{ctf}}=\Dbc\), and we write \(\Dbc\) in calculations that only concern these cases. With integral coefficients finite Tor amplitude is automatic, and membership in \(\Dbc(T,\mathcal O_E)\) can be tested modulo \(\ell\) (Proposition \bpout{3.11}).

The proof passes through finite rings \(A\) with \(\ell^rA=0\), such as \(\mathcal O_E/\ell^n\). For these, \(D^b_{\mathrm{ctf}}\) is the category of Definition \bpout{3.1}, which on the pro-étale site is identified with the pullback of its étale counterpart (Corollary \bpout{2.44}). Here the Tor condition makes derived change of coefficients preserve boundedness, as the example after Proposition \bpout{7.10} shows. When we compare coefficient rings, we write the ring as a subscript of the functor, as in \(Rf_{!,A}\) or \(f^!_{\mathcal O_E}\); otherwise we omit the subscript.

With coefficients \(A\), we have \(Rf_!=R\bar f_*\circ j_!\) for a compactification \(f=\bar f\circ j\), and on \(D^b_{\mathrm{ctf}}\) these functors and their right adjoints are those of Theorem \bpout{5.31}, restricted in Theorem \bpout{6.20}. Proper base change, the projection formula and the localization triangle are the identities of Section \bpout{6.1}, which we may prove on the étale sites by Proposition \bpout{5.15}. For an étale map, extension by zero is exact on sheaves of \(A\)-modules on the pro-étale sites; on the derived categories it is left adjoint to inverse image, and the counit is given on geometric stalks by the sum over the fibre.

Property (2) implies that the exceptional base change map is an isomorphism for smooth \(f\), because \(\gamma_f\) and \(\gamma_{f'}\) are isomorphisms.

For étale \(f\) we have \(d=0\), and the theorem says that \(f^{!}\simeq f^{*}\) for every admissible \(\Lambda\). For a smooth curve \(C\) with structure map \(a:C\to\operatorname{Spec}k\), it gives \(a^!\mathbf Z_{\ell,\operatorname{Spec}k}\simeq\mathbf Z_{\ell,C}(1)[2]\), a single sheaf in degree \(-2\).

For smooth \(f\) the whole functor \(f^!\) is determined by its value on the constant sheaf, through the map \(f^*K\otimes^{\mathbf L}f^!\Lambda_Y\to f^!K\) of Proposition proposition-exceptional-inverse-image-smooth-morphism. Smoothness is needed here, as the following example shows. Let \(i:y\to C\) be the inclusion of a rational point of a smooth curve, with open complement \(j:U\to C\), and take \(K=j_!\Lambda_U\). Then the analogous map \(i^*K\otimes^{\mathbf L}i^!\Lambda_C\to i^!K\) has source zero, since the stalk of \(j_!\Lambda_U\) at \(y\) vanishes, but the localization triangle gives \(i^!j_!\Lambda_U\simeq(i^*Rj_*\Lambda_U)[-1]\neq0\).

We give \(f^!\Lambda_Y\) a name rather than working with \(\Lambda_X(d)[2d]\) directly, because the two are not yet known to be isomorphic. Moreover, the complex \(f^!\Lambda_Y\) is attached to \(f\) without choices, and it satisfies base change (Proposition proposition-exceptional-inverse-image-smooth-morphism). Proving the theorem thus amounts to constructing an isomorphism between it and \(\Lambda_X(d)[2d]\) that does not depend on charts. In this sense the complex is the analogue of the shifted orientation sheaf \(\mathrm{or}_{X/Y}[m]\), and for the structure map of a smooth scheme it is the dualizing complex of Poincaré duality (Corollary \hyperref[bp96]{9.28} and Theorem \bpout{10.12}).

% object: definition-relative-dualizing-complex (blueprint/chap-Smooth-purity/sec-Statement/definition-relative-dualizing-complex/_section.md)
\phantomsection\label{bp11}%
\begin{bpstatement}{Definition~9.4}{ (Relative dualizing complex)}
For smooth \(f:X\to Y\) and any admissible \(\Lambda\), the \emph{relative dualizing complex} is

\[
\omega_f=f^{!}\Lambda_Y
\in D^b_{\mathrm{ctf}}(X,\Lambda).
\tag{9.3}
\]
\end{bpstatement}

% text: prose-09-06
\phantomsection\label{bp12}%
For finite coefficients, \(\omega_f\) can be computed on the pro-étale site or, equivalently, on the étale site (Proposition \bpout{5.15}).


% object: sec-The-tensor-product-map (blueprint/chap-Smooth-purity/sec-The-tensor-product-map/_section.md), content hidden
\subsection*{\texorpdfstring{9.2\quad The tensor product map}{9.2 The tensor product map}}
\phantomsection\label{bp13}%
\addcontentsline{toc}{subsection}{\texorpdfstring{9.2\quad The tensor product map}{9.2 The tensor product map}}
% text: prose-09-07
\phantomsection\label{bp14}%
With finite coefficients, smooth purity says that for \(f\) smooth of pure relative dimension \(d\) the exceptional inverse image \(f^!\) is the ordinary inverse image, twisted by \((d)\) and shifted by \([2d]\) (Proposition \hyperref[bp37]{9.12}). The first step is the map \(f^*K\otimes^{\mathbf L}\omega_f\to f^!K\), an isomorphism by Proposition proposition-exceptional-inverse-image-smooth-morphism, which reduces the problem to the single complex \(\omega_f\).

Fix \(n\geq1\). We carry out the geometric calculation with coefficients \(\mathbf Z/\ell^n\mathbf Z\), and until Section \hyperref[bp39]{9.5} we omit the subscript, so that \(f^*\), \(Rf_!\), \(f^!\), \(\omega_f\) and \(\gamma_f\) refer to this cyclic coefficient ring. Afterwards the calculation extends along \(\mathbf Z/\ell^n\mathbf Z\to A\) when \(\ell^nA=0\), and in this way Proposition \hyperref[bp37]{9.12} gives the result for every finite admissible ring \(A\).

Let \(f:X\to Y\) be smooth of pure relative dimension \(d\) and \(K\in D^b_{\mathrm{ctf}}(Y_{\mathrm{pro}\acute{e}t},\mathbf Z/\ell^n\mathbf Z)\). From the projection formula and the counit \(Rf_!\omega_{f}\to(\mathbf Z/\ell^n\mathbf Z)_{Y}\) we get \(Rf_!\bigl(f^{*}K\otimes^{\mathbf L}\omega_{f}\bigr)
\simeq K\otimes^{\mathbf L} Rf_!\omega_{f}\to K\), whose adjoint is the canonical map

\[
\mu_{f,K}:f^{*}K\otimes^{\mathbf L}_{(\mathbf Z/\ell^n\mathbf Z)_{X}}\omega_{f}
\longrightarrow f^{!}K.
\tag{9.4}
\]

To construct the orientation, we compare \(\omega\) on the fibres of a family of sections, and for this we need base change. Consider a cartesian square

\[
\begin{CD}
X' @>{g'}>> X\\
@V{f'}VV @VV{f}V\\
Y' @>{g}>> Y,
\end{CD}
\tag{9.5}
\]

with \(f\) smooth of pure relative dimension \(d\). For it we have the exceptional base change map of Definition \hyperref[bp5]{9.1}. By Proposition proposition-exceptional-inverse-image-smooth-morphism and (3), with \(\Lambda=\mathbf Z/\ell^n\mathbf Z\), the map (9.4) and the exceptional base change map

\[
g'^{*}f^{!}K\xrightarrow{\sim}(f')^{!}g^{*}K
\tag{9.6}
\]

are isomorphisms for \(K\in D^b_{\mathrm{ctf}}(Y_{\mathrm{pro}\acute{e}t},\mathbf Z/\ell^n\mathbf Z)\), and they are compatible with each other. Moreover, on a chart \(f|_U=p_d\circ e\) we have \(\omega_f|_U\simeq(\mathbf Z/\ell^n\mathbf Z)_U(d)[2d]\). Hence \(\omega_f[-2d]\) is a locally constant sheaf of free \(\mathbf Z/\ell^n\mathbf Z\)-modules of rank one.


% object: sec-The-trace-map-of-affine-space (blueprint/chap-Smooth-purity/sec-The-trace-map-of-affine-space/_section.md), content hidden
\subsection*{\texorpdfstring{9.3\quad The trace map of affine space}{9.3 The trace map of affine space}}
\phantomsection\label{bp15}%
\addcontentsline{toc}{subsection}{\texorpdfstring{9.3\quad The trace map of affine space}{9.3 The trace map of affine space}}
% text: prose-09-08
\phantomsection\label{bp16}%
Every identification of \(f^!\) in this chapter is ultimately an adjunction whose counit is integration along affine lines, so we first make that integration explicit. The affine line is not proper, so we cannot integrate on it directly; instead we integrate on \(\mathbf P^1\). There, on a geometric fibre \(H^2(\mathbf P^1,\mathbf Z/\ell^n(1))\simeq\mathbf Z/\ell^n\) is generated by \(c_1(\mathcal O(1))\), the class of a point, and the trace map of \(\mathbf P^1\) reads off the coefficient of this class. The trace map of \(\mathbf A^1\) is then its composite with the map \(H^2_c(\mathbf A^1,\mathbf Z/\ell^n(1))\to H^2(\mathbf P^1,\mathbf Z/\ell^n(1))\) given by extension by zero. This construction works over an arbitrary base \(S\) and commutes with base change, and iterating it gives the trace map (9.11) of \(\mathbf A^d\), which is the counit of the adjunction on a chart.

For a scheme \(T\), Chern classes are defined by the Kummer sequence on \(T_{\acute{e}t}\),

\[
1\longrightarrow\mu_{\ell^n}\longrightarrow\mathbf G_m
\xrightarrow{\ (\cdot)^{\ell^n}\ }\mathbf G_m
\longrightarrow1,
\tag{9.7}
\]

whose boundary \(H^1(T_{\acute{e}t},\mathbf G_m)\to
H^2(T_{\acute{e}t},(\mathbf Z/\ell^n\mathbf Z)_T(1))\) sends a line bundle \(L\) to \(c_1(L)\). For a scheme \(S\), let \(q:\mathbf P^1_S\to S\) be the projection. By Lemma \bpout{6.4}, the unit and \(c_1(\mathcal O_{\mathbf P^1_S}(1))\) give an isomorphism, compatible with base change,

\[
Rq_*(\mathbf Z/\ell^n\mathbf Z)_{\mathbf P^1_S}
\simeq(\mathbf Z/\ell^n\mathbf Z)_{S}\oplus(\mathbf Z/\ell^n\mathbf Z)_{S}(-1)[-2],
\tag{9.8}
\]

and the \emph{trace map} \(\operatorname{Tr}_q:Rq_*(\mathbf Z/\ell^n\mathbf Z)_{\mathbf P^1_S}(1)[2]
\to(\mathbf Z/\ell^n\mathbf Z)_{S}\) is the projection onto the second summand, after twisting. It sends \(c_1(\mathcal O_{\mathbf P^1_{\bar s}}(1))\) on each geometric fibre to \(1\).

Let \(p_1,p_2\) be the projections of \(\mathbf P^1_S\times_S\mathbf P^1_S\), and let \(x_0,x_1\) and \(y_0,y_1\) be the pullbacks along \(p_1\) and \(p_2\) of the homogeneous coordinates of \(\mathbf P^1_S\). The diagonal \(\Delta\) is the zero locus of the section \(x_0y_1-x_1y_0\) of \(p_1^*\mathcal O(1)\otimes p_2^*\mathcal O(1)\), so that \(\mathcal O(\Delta)\) is this line bundle and \(x_0y_1-x_1y_0\) is its canonical section. Lemma \bpout{6.3} provides the cycle class of \(\Delta\),

\[
\delta:\Delta_*(\mathbf Z/\ell^n\mathbf Z)_{\mathbf P^1_S}
\longrightarrow(\mathbf Z/\ell^n\mathbf Z)_{\mathbf P^1_S\times_S\mathbf P^1_S}(1)[2],
\]

the Kummer boundary of \(\mathcal O(\Delta)\) with its canonical section, compatible with every base change on \(S\). Let \(\operatorname{Tr}_{p_i}\) be the pullback of \(\operatorname{Tr}_q\). Then the diagonal integrates to \(1\) along either factor, that is,

\[
\operatorname{Tr}_{p_i}\circ Rp_{i*}\delta
=\operatorname{id}_{(\mathbf Z/\ell^n\mathbf Z)_{\mathbf P^1_S}}
\qquad(i=1,2).
\tag{9.9}
\]

This is (6.6) of Proposition \bpout{6.7}. To see it directly, note that both sides are endomorphisms of a constant sheaf, so it suffices to check (9.9) on geometric fibres of \(p_i\), where the diagonal is a single rational point with line bundle \(\mathcal O(1)\), and \(\delta\) restricts to \(c_1(\mathcal O(1))\), which \(\operatorname{Tr}_q\) sends to \(1\).

For \(\Lambda=\mathbf Z/\ell^n\mathbf Z\), Proposition \bpout{6.7} gives the adjunction \(Rq_*\dashv q^*(-)(1)[2]\), and its unit and counit, (6.7) and (6.8) there, take the following form. The counit is integration along the fibres,

\[
\varepsilon_K:Rq_*\bigl(q^{*}K(1)[2]\bigr)
\simeq K\otimes^{\mathbf L} Rq_*(\mathbf Z/\ell^n\mathbf Z)_{\mathbf P^1_S}(1)[2]
\xrightarrow{\ \operatorname{id}_K\otimes\operatorname{Tr}_q\ } K,
\tag{9.10}
\]

and the unit is the class of the diagonal,

\[
\eta_M:\ M\simeq Rp_{2*}\Delta_*M
\xrightarrow{\ Rp_{2*}(\delta\otimes p_1^*M)\ } Rp_{2*}p_{1}^{*}M(1)[2]
\simeq q^{*}Rq_*M(1)[2],
\]

using the projection formulas and proper base change. Both commute with base change on \(S\), and the triangle identities reduce to (9.9).

Let \(j:\mathbf A^1_S\hookrightarrow\mathbf P^1_S\) be the complement of the section at infinity, and let \(p=q\circ j\). Composing this adjunction with \(j_!\dashv j^*\) we obtain an adjunction \(Rp_!\dashv p^*(-)(1)[2]\) between \(D^b_{\mathrm{ctf}}(\mathbf A^1_{S,\mathrm{pro}\acute{e}t},\mathbf Z/\ell^n\mathbf Z)\) and \(D^b_{\mathrm{ctf}}(S_{\mathrm{pro}\acute{e}t},\mathbf Z/\ell^n\mathbf Z)\), whose counit is the counit for the open immersion followed by \(\varepsilon_K\). At the constant sheaf this counit is the trace map of the affine line \(\operatorname{Tr}_p:Rp_!(\mathbf Z/\ell^n\mathbf Z)_{\mathbf A^1_S}(1)[2]
\to(\mathbf Z/\ell^n\mathbf Z)_S\) of (9.1), and it is an isomorphism. Indeed, we apply \(Rq_*\) to the localization triangle for the section at infinity. In the resulting triangle, restriction to infinity is the identity on the first summand of (9.8), and it kills the second summand, because the pullback of \(\mathcal O(1)\) along the section at infinity is the trivial line bundle on \(S\), and so the restriction of \(c_1(\mathcal O(1))\) is \(0\). Hence \(Rp_!(\mathbf Z/\ell^n\mathbf Z)_{\mathbf A^1_S}\simeq
(\mathbf Z/\ell^n\mathbf Z)_S(-1)[-2]\) is the second summand. Over a geometric point this is \(H^2_c(\mathbf A^1,(\mathbf Z/\ell^n\mathbf Z)(1))\simeq
\mathbf Z/\ell^n\mathbf Z\), generated by the class \(\kappa\) that maps to the class of a point in \(H^2(\mathbf P^1,(\mathbf Z/\ell^n\mathbf Z)(1))\).

For affine space, let \(p_j:\mathbf A^j_S\to S\) be the projection and let \(r_j:\mathbf A^j_S\to\mathbf A^{j-1}_S\) be the map forgetting the last coordinate. We define \(\operatorname{Tr}_{p_j}\) inductively as the composite

\[
\begin{aligned}
Rp_{j!}(\mathbf Z/\ell^n\mathbf Z)_{\mathbf A^j_S}(j)[2j]
&\simeq Rp_{j-1!}\bigl(Rr_{j!}(\mathbf Z/\ell^n\mathbf Z)_{\mathbf A^j_S}(1)[2]\bigr)(j-1)[2j-2]\\
&\xrightarrow{\ \operatorname{Tr}_{p_{j-1}}\circ Rp_{j-1!}(\operatorname{Tr}_{r_j})\ }
(\mathbf Z/\ell^n\mathbf Z)_S,
\end{aligned}
\]

using \(Rp_{j!}=Rp_{j-1!}Rr_{j!}\) and the projection formula. For \(j=d\) this is the \emph{trace map of affine space}, an isomorphism (the identity when \(d=0\))

\[
\operatorname{Tr}_{p_d}:Rp_{d!}(\mathbf Z/\ell^n\mathbf Z)_{\mathbf A^d_S}(d)[2d]
\xrightarrow{\sim}(\mathbf Z/\ell^n\mathbf Z)_{S}.
\tag{9.11}
\]

By Lemma lemma-affine-space-etale-localization, \(Rp_{d!}\dashv
p_d^*(-)(d)[2d]\), and the counit is the composite of the \(d\) counits for the affine lines, which at the constant sheaf is (9.11). Like those of the affine line, the unit and counit commute with base change. On a geometric fibre, the external product \(\kappa\boxtimes\cdots\boxtimes\kappa\in
H_c^{2d}(\mathbf A^d_{\bar s},(\mathbf Z/\ell^n\mathbf Z)(d))\) integrates to \(1\), since by the projection formula, integrating along the last coordinate replaces the last factor by its integral, which is \(1\).

Objects of a derived category do not glue in general. For example, suppose \(X=U\cup V\) and we know objects on \(U\) and on \(V\) isomorphic to \(f^!K|_U\) and \(f^!K|_V\). Then the two isomorphisms restricted to \(U\cap V\) may differ by an automorphism, and nothing tells us which identification to glue along. On a chart \(U\to\mathbf A^d_Y\) we have an explicit candidate for \(f^!K|_U\), namely \(e^*p_d^*K(d)[2d]\), and we want to compare such candidates with \(f^!K\) and later, for towers, to glue \(f^!\) from them. For this we keep the counit together with the object. The point is that a pair consisting of an object and a counit that satisfies a universal property is unique up to a unique isomorphism compatible with the counit. Hence local pairs can be compared on overlaps without any choice of transition maps. We use this for the comparison following Definition \hyperref[bp17]{9.5} and in the construction of the right adjoint for towers (Proposition \hyperref[bp58]{9.18}).

% object: definition-representing-objects (blueprint/chap-Smooth-purity/sec-The-trace-map-of-affine-space/definition-representing-objects/_section.md)
\phantomsection\label{bp17}%
\begin{bpstatement}{Definition~9.5}{ (Representing objects)}
Let \(f\) be smooth, and let us work in \(D(Y,\mathbf Z/\ell^n\mathbf Z)\) and \(D(X,\mathbf Z/\ell^n\mathbf Z)\) on either the étale or the pro-étale sites. Let \(K\) be an object on \(Y\). We say that a pair \((H,\varepsilon)\), where \(H\) is an object on \(X\) and \(\varepsilon:Rf_!H\to K\) is a morphism, \emph{represents} \(\operatorname{Hom}(Rf_!(-),K)\) if \(b\mapsto\varepsilon\circ Rf_!(b)\) is a bijection \(\operatorname{Hom}(M,H)\to\operatorname{Hom}(Rf_!M,K)\) for all \(M\in D^b_{\mathrm{ctf}}(X,\mathbf Z/\ell^n\mathbf Z)\).
\end{bpstatement}

% text: prose-09-09
\phantomsection\label{bp18}%
The geometric computations with finite coefficients in Section \bpout{6.1} are stated on the pro-étale site and proved by calculations on the étale site. By Corollary \bpout{2.44}, Theorem \bpout{2.38} with (2.10), and Proposition \bpout{5.15}, the pro-étale and étale categories \(D^b_{\mathrm{ctf}}(-,A)\), and the localizing subcategories they generate, are identified compatibly with pullback, direct image and extension by zero. In particular \(Rf_!=R\bar f_*\circ j_!\) is computed on the étale sites. Consequently, an adjunction between \(Rf_!\) and some functor \(G\) on the étale categories \(D^b_{\mathrm{ctf}}(-_{\acute{e}t},A)\) transfers to the pro-étale categories \(D^b_{\mathrm{ctf}}(-_{\mathrm{pro}\acute{e}t},A)\). Indeed, the Hom groups agree by full faithfulness, and the two functors \(Rf_!\) correspond to each other.

By Proposition \bpout{5.15}, the functor \(f^!\) on \(\Loc(D^b_{\mathrm{ctf}}(-,\mathbf Z/\ell^n\mathbf Z))\) of Theorem \bpout{5.31} is the étale functor \(f^!\) on \(D(-_{\acute{e}t},\mathbf Z/\ell^n\mathbf Z)\) of Section \bpout{6.1}, and it preserves \(D^b_{\mathrm{ctf}}\) by Theorem \bpout{6.20}. Hence, for \(f:X\to Y\) smooth of pure relative dimension \(d\), Proposition proposition-exceptional-inverse-image-smooth-morphism gives the local form of \(Rf_!\dashv f^!\): on an open \(U\subset X\) with \(f|_U=p_d\circ e\), where \(e:U\to\mathbf A^d_Y\) is étale,

\[
f^!K|_U\simeq e^*p_{d}^*K(d)[2d],
\]

with counit the sum over the fibres of \(e\) followed by (9.11) tensored with \(K\).

Since a right adjoint is unique up to a unique isomorphism preserving the counit, these counits on charts agree on overlaps with the counit of \(Rf_!\dashv f^!\). With finite coefficients, then, \(f^!\) exists first and the charts describe it. For towers, by contrast, there is no global functor yet, and we glue \(f^!\) from models on charts by a Mayer–Vietoris argument, which uses the triangle for an open cover \(X=U\cup V\) and \(N\in D^b_{\mathrm{ctf}}(X_{\mathrm{pro}\acute{e}t},\mathbf Z/\ell^n\mathbf Z)\):

\[
j_{U\cap V!}(N|_{U\cap V})
\xrightarrow{\ (1,-1)\ }
j_{U!}(N|_U)\oplus j_{V!}(N|_V)
\xrightarrow{\ (1,1)\ }N\longrightarrow{} .
\tag{9.12}
\]

Suppose that \(H\in D^b_{\mathrm{ctf}}(X_{\mathrm{pro}\acute{e}t},\mathbf Z/\ell^n\mathbf Z)\) and \(\varepsilon:Rf_!H\to K\) restrict to pairs that represent \(\operatorname{Hom}(Rf_!(-),K)\) (Definition \hyperref[bp17]{9.5}) on \(U\), \(V\) and \(U\cap V\), where the restriction of \(\varepsilon\) to \(U\) is the composite \(R(f|_U)_!(H|_U)=Rf_!j_{U!}(H|_U)\to Rf_!H\to K\). For \(M\in D^b_{\mathrm{ctf}}(X_{\mathrm{pro}\acute{e}t},\mathbf Z/\ell^n\mathbf Z)\) we apply \(\operatorname{Hom}(-,H)\) and \(\operatorname{Hom}(Rf_!(-),K)\) to (9.12) with \(N=M\). The maps \(b\mapsto\varepsilon\circ Rf_!(b)\) form a morphism between the two long exact sequences, which is bijective on the terms coming from \(U\), \(V\) and \(U\cap V\). By the five lemma it is bijective on \(\operatorname{Hom}(M,H)\), which means that \((H,\varepsilon)\) represents \(\operatorname{Hom}(Rf_!(-),K)\). Since such a pair is unique up to a unique isomorphism compatible with \(\varepsilon\), we get \(H\simeq f^!K\) compatibly with the counits. In Section \hyperref[bp39]{9.5} we use this gluing to realize the compatible counits with finite coefficients as a single map of towers.

On a vector bundle, the choice of a local trivialization transports (9.11) to a trace map. These maps do not depend on the trivialization, so they glue to a global map. This is the model case of the orientation below, because deformation to the normal bundle will reduce the dualizing complex of a smooth morphism to that of a vector bundle.

% object: proposition-purity-vector-bundles (blueprint/chap-Smooth-purity/sec-The-trace-map-of-affine-space/proposition-purity-vector-bundles/_section.md)
\phantomsection\label{bp19}%
\begin{bpstatementit}{Proposition~9.6}{ (Purity for vector bundles)}
Let \(\pi:V\to S\) be a vector bundle of rank \(d\). For \(K\in D^b_{\mathrm{ctf}}(S_{\mathrm{pro}\acute{e}t},\mathbf Z/\ell^n\mathbf Z)\), the trace map induces an isomorphism

\[
\pi^{*}K(d)[2d]\xrightarrow{\sim}\pi^{!}K.
\]
\end{bpstatementit}

% proof: proof/proposition-purity-vector-bundles
\begin{proof}[{Proof}]\phantomsection\label{bp20}%
We transport the trace map of affine space along local trivializations, show that the result does not depend on the trivialization, so that the local maps glue to a trace isomorphism \(\operatorname{Tr}_\pi\) for \(V\), and take the map \(\pi^*K(d)[2d]\to\pi^!K\) to be its adjoint.

\emph{Local trace maps.} Locally on \(S\) the bundle is affine space, where the trace map has already been constructed. Over an open \(U\subset S\) on which \(V\) is trivial, a trivialization \(\tau:V\times_S U\xrightarrow{\sim}\mathbf A^d_U\) transports (9.11) to an isomorphism

\[
\operatorname{Tr}_{\tau}:
R(\pi|_{V\times_S U})_!(\mathbf Z/\ell^n\mathbf Z)_{V\times_S U}(d)[2d]
\xrightarrow{\sim}(\mathbf Z/\ell^n\mathbf Z)_U.
\]

These maps commute with base change in \(U\), because the construction uses only \(c_1(\mathcal O_{\mathbf P^1_U}(1))\), the splitting (9.8) and localization at infinity: pullback preserves the Kummer class, and proper base change and localization identify the pullbacks of the constructions with those over the new base.

\emph{Independence of the trivialization.} A change of trivialization multiplies the trace map by a locally constant unit on \(\GL_d\), and this unit is \(1\). As two trivializations differ by a matrix with values in \(\GL_d\), it suffices to study one universal change of coordinates, parametrized by \(\GL_{d,S}\) itself, and then pull back. Let

\[
\theta:\mathbf A^d_{\GL_{d,S}}\longrightarrow\mathbf A^d_{\GL_{d,S}},
\qquad (g,v)\longmapsto(g,gv),
\]

be the universal linear change of coordinates, an automorphism over \(\GL_{d,S}\). It induces an automorphism of the compactly supported direct image of \((\mathbf Z/\ell^n\mathbf Z)_{\mathbf A^d_{\GL_{d,S}}}(d)[2d]\) to \(\GL_{d,S}\), which the trace map identifies with the constant sheaf \((\mathbf Z/\ell^n\mathbf Z)_{\GL_{d,S}}\) in degree zero, so that this automorphism is multiplication by a section

\[
u\in H^0\bigl((\GL_{d,S})_{\acute{e}t},
(\mathbf Z/\ell^n\mathbf Z)_{\GL_{d,S}}^{\times}\bigr),
\]

a locally constant function with values in \((\mathbf Z/\ell^n\mathbf Z)^{\times}\). It is therefore constant on each geometric fibre \(\GL_{d,\bar s}\), which is connected: for \(d\geq1\) it is a nonempty open subset of the affine space of \(d\times d\) matrices, and for \(d=0\) it is a point. As \(\theta\) is the identity at the identity matrix, which lies on every geometric fibre, \(u=1\).

The local trace maps agree on overlaps of trivializing opens. Indeed, let \(\tau\) and \(\tau'\) be two trivializations over \(U\). Then \(\tau'\circ\tau^{-1}\) is the pullback of \(\theta\) along the transition matrix \(U\to\GL_{d,S}\), so by base change the equality \(u=1\) implies \(\operatorname{Tr}_{\tau'}=\operatorname{Tr}_{\tau}\). Moreover, morphisms between source and target form a sheaf, as both are locally sheaves in degree zero. Hence the local maps glue over a trivializing cover to an isomorphism

\[
\operatorname{Tr}_{\pi}:
R\pi_!(\mathbf Z/\ell^n\mathbf Z)_V(d)[2d]
\xrightarrow{\sim}(\mathbf Z/\ell^n\mathbf Z)_S.
\]

\emph{The map \(\pi^*K(d)[2d]\to\pi^!K\).} For general \(K\), the projection formula and \(\operatorname{Tr}_\pi\) give

\[
R\pi_!\bigl(\pi^{*}K(d)[2d]\bigr)
\simeq K\otimes^{\mathbf L} R\pi_!(\mathbf Z/\ell^n\mathbf Z)_V(d)[2d]
\xrightarrow{\ \mathrm{id}_K\otimes\operatorname{Tr}_{\pi}\ }K.
\]

Its adjoint under \(R\pi_!\dashv\pi^{!}\) (Proposition proposition-exceptional-inverse-image-smooth-morphism) is the asserted map. We show that it is an isomorphism. This may be checked locally on \(V\), because \(\pi^!\) commutes with restriction to opens. Over each trivializing open, \(\operatorname{Tr}_\pi\) is \(\operatorname{Tr}_\tau\) by construction, so there the map is the isomorphism given by the adjunction for affine space following (9.11).
\end{proof}

% text: prose-09-10
\phantomsection\label{bp21}%
The trace maps of vector bundles commute with base change on \(S\), because the local trace maps do.


% object: sec-Finite-coefficients-3 (blueprint/chap-Smooth-purity/sec-Finite-coefficients-3/_section.md), content hidden
\subsection*{\texorpdfstring{9.4\quad Finite coefficients}{9.4 Finite coefficients}}
\phantomsection\label{bp22}%
\addcontentsline{toc}{subsection}{\texorpdfstring{9.4\quad Finite coefficients}{9.4 Finite coefficients}}
% text: prose-09-11
\phantomsection\label{bp23}%
Since \(f^!K\simeq f^*K\otimes^{\mathbf L}\omega_f\), smooth purity with finite coefficients comes down to a single global isomorphism \(\omega_f\simeq(\mathbf Z/\ell^n\mathbf Z)_X(d)[2d]\). We obtain it from the diagonal \(\Delta:X\to X\times_YX\), which is a section of the smooth projection \(p_1\). Deformation to the normal bundle moves \(\Delta\) to the zero section of the relative tangent bundle \(T_{X/Y}\), where Proposition \hyperref[bp19]{9.6} gives the isomorphism.

We use the diagonal because it is attached to \(f\) without choices, and because \(\omega_f\) is recovered from it: \(\omega_f\simeq\Delta^*\omega_{p_1}\) by base change. Near the diagonal, however, \(X\times_YX\) looks like a neighbourhood of the zero section in \(T_{X/Y}\) only to first order, and the deformation makes this precise without choosing a tubular neighbourhood. The comparison it yields, Lemma \hyperref[bp27]{9.9}, is determined by the requirement that it be the identity at \(t=1\).

The deformation argument works for any section of a smooth morphism. Let \(g:V\to S\) be smooth of pure relative dimension \(d\) and let \(s:S\to V\) be a section of \(g\) that is a closed immersion. Then \(s\) is a regular immersion of codimension \(d\), and we let \(N_s\) be its normal bundle, with projection \(\pi:N_s\to S\) and zero section \(z:S\to N_s\).

% object: definition-deformation-normal-bundle (blueprint/chap-Smooth-purity/sec-Finite-coefficients-3/definition-deformation-normal-bundle/_section.md)
\phantomsection\label{bp24}%
\begin{bpstatement}{Definition~9.7}{ (Deformation to the normal bundle)}
The \emph{deformation to the normal bundle} of \(s\) is the open subscheme \(\mathcal D_{s(S)}V\) of the blowup of \(V\times\mathbf A^1\) along \(s(S)\times\{0\}\) obtained by removing the strict transform of \(V\times\{0\}\).
\end{bpstatement}

% text: prose-09-12
\phantomsection\label{bp25}%
It comes with a smooth morphism \(\widetilde g:\mathcal D_{s(S)}V\to\mathbf A^1_S\) and a section \(\widetilde s:\mathbf A^1_S\to\mathcal D_{s(S)}V\) that is a closed immersion. Over \(t=1\) the pair \((\widetilde g,\widetilde s)\) restricts to \((g,s)\), and over \(t=0\) to \((\pi,z)\). To see this in local terms, let \(t\) be the coordinate of \(\mathbf A^1\), and let \(x_1,\ldots,x_d\) be functions on an open neighbourhood \(U\) of \(s(S)\) in \(V\) which are étale coordinates over \(S\) and generate the ideal of \(s(S)\). On the part of \(\mathcal D_{s(S)}V\) lying over \(U\times\mathbf A^1\), the quotients \(u_i=x_i/t\) are regular functions, so that \(x_i=tu_i\) there. At \(t=0\), then, the function \(u_i\) restricts to the linear function on \(N_s\) given by the class of \(x_i\) in the conormal sheaf. Since these classes form a basis of the conormal sheaf, the \(u_i\) restrict to linear coordinates on \(N_s\).

% object: example-deformation-origin-line (blueprint/chap-Smooth-purity/sec-Finite-coefficients-3/example-deformation-origin-line/_section.md)
\phantomsection\label{bp26}%
\begin{bpstatement}{Example~9.8}{ (The origin in the affine line)}
Let \(S=\operatorname{Spec}k\), \(V=\mathbf A^1_k\) with coordinate \(x\), and let \(s\) be the origin. The blowup of \(V\times\mathbf A^1=\mathbf A^2\), with coordinates \((x,t)\), at the origin is covered by two charts. On the chart with coordinates \((u,t)\) and \(x=tu\) the exceptional divisor is \(t=0\); on the chart with coordinates \((x,v)\) and \(t=xv\) it is \(x=0\), and the strict transform of the line \(V\times\{0\}=\{t=0\}\) is \(v=0\). Since a point of the second chart with \(v\neq0\) lies on the first chart, with \(u=1/v\), removing the strict transform leaves exactly the first chart:

\[
\mathcal D_{s(S)}V=\operatorname{Spec}k[u,t]\longrightarrow
V\times\mathbf A^1,\qquad (u,t)\longmapsto(tu,t),
\]

with \(\widetilde g\) the projection to \(t\) and \(\widetilde s\) the line \(u=0\). For \(t=c\neq0\) the fibre is the line \(V\) itself, with coordinate \(u=x/c\) and the origin as section. For \(t=0\) the fibre is the \(u\)-line, mapped to the origin of \(V\); it is the tangent line \(N_s=T_0\mathbf A^1\), and \(u\) is the linear coordinate given by the class of \(x\) in \(\mathfrak m/\mathfrak m^2\). So the family rescales the line by \(1/t\) and in the limit keeps only the first-order information at the section. Here the family is the product \(\mathbf A^1_u\times\mathbf A^1_t\), so \(\widetilde s^*\omega_{\widetilde g}\) is visibly constant along \(t\). In general this holds only étale locally, after a choice of coordinates, and the point of Lemma \hyperref[bp27]{9.9} is that the comparison normalized at \(t=1\) does not depend on that choice.
\end{bpstatement}

% object: lemma-section-normal-bundle (blueprint/chap-Smooth-purity/sec-Finite-coefficients-3/lemma-section-normal-bundle/_section.md)
\phantomsection\label{bp27}%
\begin{bpstatementit}{Lemma~9.9}{ (A section and its normal bundle)}
There is a canonical isomorphism

\[
s^{*}\omega_{g}\xrightarrow{\sim}z^{*}\omega_{\pi}.
\tag{9.13}
\]

It is obtained by restricting to \(t=0\) the unique isomorphism \(\widetilde s^{*}\omega_{\widetilde g}\to a^{*}s^{*}\omega_{g}\) that is the identity at \(t=1\), and inverting. Here \(a:\mathbf A^1_S\to S\) is the projection.
\end{bpstatementit}

% proof: proof/lemma-section-normal-bundle
\begin{proof}[{Proof}]\phantomsection\label{bp28}%
The fibres of \(\widetilde s^{*}\omega_{\widetilde g}\) at \(t=1\) and \(t=0\) are the two sides of (9.13), so it suffices to show that this shifted sheaf of rank one is pulled back from \(S\), étale locally on \(S\). In positive characteristic a locally constant sheaf on \(\mathbf A^1_S\) need not be pulled back from \(S\), so we use coordinates. Étale locally on \(S\), near the section the deformation is étale over \(\mathbf A^d_{\mathbf A^1_S}=\mathbf A^d_S\times\mathbf A^1\), a family pulled back from \(S\), and hence so is \(\widetilde s^{*}\omega_{\widetilde g}[-2d]\). We then normalize the local comparisons at \(t=1\); since the geometric fibres of \(\mathbf A^1_S\to S\) are connected, this makes them agree on overlaps. Finally, they glue because all objects involved are shifted locally constant sheaves.

\emph{Step 1: the objects.} The restrictions of \(\widetilde s^{*}\omega_{\widetilde g}\) to \(t=1\) and \(t=0\) are \(s^{*}\omega_{g}\) and \(z^{*}\omega_{\pi}\), by Proposition proposition-exceptional-inverse-image-smooth-morphism applied to the inclusions of these fibres. Proposition proposition-exceptional-inverse-image-smooth-morphism shows that both \(\widetilde s^{*}\omega_{\widetilde g}[-2d]\) and \(s^{*}\omega_{g}[-2d]\) are locally constant sheaves, free of rank one over \(\mathbf Z/\ell^n\mathbf Z\). Hence morphisms between two such complexes are morphisms of sheaves, and they glue.

\emph{Step 2: local construction.} Étale locally on \(S\) we reduce to the zero section of affine space, whose deformation is pulled back from \(S\). After passing to an étale cover \(\{S_\alpha\to S\}_\alpha\) we may choose, for the base change of the section \(s_\alpha:S_\alpha\to V\times_SS_\alpha\), a neighbourhood \(U_\alpha\) of \(s_\alpha(S_\alpha)\) and an étale map \(e_\alpha:U_\alpha\to\mathbf A^d_{S_\alpha}\) sending \(s_\alpha\) to the zero section. We may assume that the inverse image of the zero section under \(e_\alpha\) is exactly \(s_\alpha(S_\alpha)\). Indeed, this inverse image is separated and étale over \(S_\alpha\), so \(s_\alpha(S_\alpha)\) is open and closed in it, and we remove the complement, which is closed in \(U_\alpha\). Then \(e_\alpha\) is an étale map of pairs whose square of centers is cartesian. Since blowing up and removing the strict transform commute with étale base change, \(e_\alpha\) induces an étale map from a neighbourhood of the deformed section in \(\mathcal D_{s(S)}V\) to the deformation space of the zero section of \(\mathbf A^d_{S_\alpha}\). The latter is \(\mathbf A^d_{\mathbf A^1_{S_\alpha}}\), with coordinates \(u_1,\ldots,u_d\) and \(t\), and its map to \(\mathbf A^d_{S_\alpha}\times\mathbf A^1\) sends \((u,t)\) to \((tu,t)\), that is, \(x_i=tu_i\), and the deformed zero section is \(u=0\).

The adjunction for affine space identifies the relative dualizing complex of the projection \(\mathbf A^d_{\mathbf A^1_{S_\alpha}}\to\mathbf A^1_{S_\alpha}\) with the constant sheaf \((\mathbf Z/\ell^n\mathbf Z)(d)[2d]\). Pulling back along the étale map just constructed, and composing adjunctions for it, gives the same identification for \(\omega_{\widetilde g}\) near the deformed section. Similarly, the chart \(e_\alpha\) gives the analogous identification of \(s^{*}\omega_{g}\) over \(S_\alpha\). Composing the two identifications we obtain an isomorphism over \(\mathbf A^1_{S_\alpha}\)

\[
\widetilde s^{*}\omega_{\widetilde g}\xrightarrow{\sim}
a^{*}s^{*}\omega_{g}.
\]

At \(t=1\) it restricts to an automorphism of \(s^{*}\omega_{g}\) over \(S_\alpha\). Composing with the pullback along \(a\) of the inverse of this automorphism, we may assume that the local isomorphism is the identity at \(t=1\).

\emph{Step 3: gluing.} The local isomorphisms from different charts agree on overlaps, because of the normalization at \(t=1\). Over \(S_\alpha\times_SS_\beta\) the two local isomorphisms that are the identity at \(t=1\) differ by an automorphism of the locally constant sheaf \(a^{*}s^{*}\omega_g[-2d]\) of rank one, that is, by a section of a locally constant sheaf of units of \(\mathbf Z/\ell^n\mathbf Z\). The locus where this section equals \(1\) is open and closed, and it contains the point \(t=1\) of each geometric fibre \(\mathbf A^1_{\bar b}\). As these fibres are connected, the section is \(1\) at every geometric point. So the local isomorphisms glue to a unique isomorphism over \(\mathbf A^1_S\) that is the identity at \(t=1\). Restricting it to \(t=0\) gives an isomorphism \(z^{*}\omega_{\pi}\simeq s^{*}\omega_g\), whose inverse is (9.13).
\end{proof}

% text: prose-09-13
\phantomsection\label{bp29}%
The orientation of \(\omega_f\) will be the comparison (9.13) for the diagonal. Two things have to be checked about it: on a coordinate chart it must be the identification given by coordinates, and it must commute with base change, as property (2) of Theorem \hyperref[bp9]{9.3} requires. Since a chart is an étale map that preserves the section, both checks need to transport (9.13) along base change on \(S\) and along étale maps that preserve the section. In both cases we use the uniqueness in Lemma \hyperref[bp27]{9.9}: any construction that is the identity at \(t=1\) produces the comparison of that lemma.

% object: lemma-naturality-under-base-change-etale (blueprint/chap-Smooth-purity/sec-Finite-coefficients-3/lemma-naturality-under-base-change-etale/_section.md)
\phantomsection\label{bp30}%
\begin{bpstatementit}{Lemma~9.10}{ (Naturality under base change and étale maps)}
Write \(\alpha_{g,s}:s^*\omega_g\xrightarrow{\sim}z^*\omega_\pi\) for the comparison of Lemma \hyperref[bp27]{9.9}.

\begin{enumerate}
\item Let \(h:S'\to S\) be any morphism for which the schemes below satisfy the standing hypotheses. Put \(V'=V\times_SS'\), and denote the base changes of the maps by \(g',s',\pi',z'\). There are canonical identifications \[
   N_{s'}\simeq N_s\times_SS',\qquad
   \mathcal D_{s'(S')}V'
   \simeq\mathcal D_{s(S)}V\times_SS'.
   \] Under these identifications and the exceptional base change isomorphisms of Proposition proposition-exceptional-inverse-image-smooth-morphism, the following square commutes: \[
   \begin{CD}
   h^*s^*\omega_g @>{h^*\alpha_{g,s}}>> h^*z^*\omega_\pi\\
   @V{\sim}VV @VV{\sim}V\\
   (s')^*\omega_{g'} @>{\alpha_{g',s'}}>> (z')^*\omega_{\pi'}.
   \end{CD}
   \]
\item Let \(e:V'\to V\) be an étale \(S\)-morphism and let \(s':S\to V'\) be a closed section such that \(es'=s\) and \(e^{-1}(s(S))=s'(S)\) scheme-theoretically. Put \(g'=ge\). The induced map \(N_{s'}\to N_s\) is a linear isomorphism, and the induced map of deformation spaces is étale. The identifications \[
   (s')^*\omega_{g'}\simeq s^*\omega_g,\qquad
   (z')^*\omega_{\pi'}\simeq z^*\omega_\pi
   \] obtained by composing adjunctions carry \(\alpha_{g',s'}\) to \(\alpha_{g,s}\).
\end{enumerate}

These comparisons respect identities, composition and interchange of base change with an étale map of pairs.
\end{bpstatementit}

% text: prose-09-14
\phantomsection\label{bp31}%
In part (1) the morphism \(h\) is arbitrary. Blowups in general commute only with flat base change, but by its description in local coordinates, the deformation space of a section of a smooth morphism commutes with arbitrary base change.

% proof: proof/lemma-naturality-under-base-change-etale
\begin{proof}[{Proof of Lemma 9.10}]\phantomsection\label{bp32}%
The geometric identifications follow from the coordinate description of the deformation space. Every comparison of dualizing complexes is then the identity at \(t=1\), and so the uniqueness in Lemma \hyperref[bp27]{9.9} identifies it with the comparison of that lemma.

\emph{The geometry.} Since \(g\) is smooth and \(s\) is a section, \(N_s\simeq s^*T_{V/S}\). Since formation of relative differentials for a smooth map commutes with arbitrary base change, and an étale map induces an isomorphism on relative tangent bundles, this proves both assertions about normal bundles.

For deformation spaces, we use the cartesian étale coordinate chart \(U\to\mathbf A^d_S\) of the proof of Lemma \hyperref[bp27]{9.9}. There the deformation is

\[
\mathcal D_{s(S)}U
\simeq U\times_{\mathbf A^d_S}\mathbf A^d_{\mathbf A^1_S},
\qquad x_i=t u_i,
\]

because this is the chart of the blowup along \((t,x_1,\ldots,x_d)\) on which the pullback of the ideal is generated by \(t\), and this description commutes with every change of base. Away from the centre, the deformation is \((V\setminus s(S))\times\mathbf G_m\), which also commutes with base change. After passing to an étale cover of \(S\) if necessary, these opens cover the deformation space, and their identifications glue to the canonical one. For an étale map of pairs as in (2), we obtain the étale map of deformation spaces by flat base change for the blowup and for the removal of the strict transform, and the coordinate description identifies its special fibre with the map of normal bundles.

\emph{The comparison of dualizing complexes.} In both parts the strategy is the same: transport the comparison over \(\mathbf A^1_S\) along the geometric identification, check that the result is still the identity at \(t=1\), and invoke uniqueness. Write

\[
\chi_{g,s}:\widetilde s^*\omega_{\widetilde g}\xrightarrow{\sim}a^*s^*\omega_g
\]

for the isomorphism of Lemma \hyperref[bp27]{9.9} that is the identity at \(t=1\), so that \(\alpha_{g,s}\) is the inverse of its restriction to \(t=0\). For (1), we apply Proposition proposition-exceptional-inverse-image-smooth-morphism to \(g\), \(\pi\) and \(\widetilde g\). These exceptional base change isomorphisms commute with restriction to \(t=1\) and \(t=0\), because they are adjoint transposes of base change for compactly supported direct image applied to the same counit, and this base change respects composition of cartesian squares. Transported through these isomorphisms, the pullback of \(\chi_{g,s}\) is therefore a comparison for \((g',s')\) that is the identity at \(t=1\), and by uniqueness it equals \(\chi_{g',s'}\). Restricting to \(t=0\) and inverting proves (1).

For (2), composing \(Re_!\dashv e^*\) with \(Rg_!\dashv g^!\) gives \(e^*\omega_g\simeq\omega_{g'}\), and likewise for the induced étale map of deformation spaces. On the fibres \(t=1\) and \(t=0\) these identifications are the ones for the original pairs and for the normal bundles, since the counit of a composite adjunction is the composite of the counits. They transport \(\chi_{g,s}\) to a comparison that is the identity at \(t=1\), which is \(\chi_{g',s'}\) by uniqueness.

For the coherence assertions we use the same uniqueness once more. The geometric identifications respect composition, and so do the identifications obtained from adjunctions, as we see by transposing composite counits. Therefore each composite in the coherence assertions is a comparison that is the identity at \(t=1\), and by uniqueness any two of them agree.
\end{proof}

% text: prose-09-15
\phantomsection\label{bp33}%
We apply Lemma \hyperref[bp27]{9.9} to the diagonal. Let \(p_1,p_2:X\times_YX\to X\) be the projections. Since \(f\) is separated, \(\Delta\) is a closed immersion, and it is a section of the smooth morphism \(p_1\). Its conormal sheaf is \(\Omega^1_{X/Y}\), so its normal bundle is the relative tangent bundle \(T_{X/Y}\), and we write \(\pi:T_{X/Y}\to X\) for its projection and \(z:X\to T_{X/Y}\) for its zero section.

% object: proposition-relative-dualizing-complex (blueprint/chap-Smooth-purity/sec-Finite-coefficients-3/proposition-relative-dualizing-complex/_section.md)
\phantomsection\label{bp34}%
\begin{bpstatementit}{Proposition~9.11}{ (The relative dualizing complex)}
Let \(f:X\to Y\) be smooth of pure relative dimension \(d\). There is a canonical isomorphism

\[
\omega_{f}\xrightarrow{\sim}(\mathbf Z/\ell^n\mathbf Z)_{X}(d)[2d].
\tag{9.14}
\]
\end{bpstatementit}

% proof: proof/proposition-relative-dualizing-complex
\begin{proof}[{Proof}]\phantomsection\label{bp35}%
We define the isomorphism as a chain of three steps: restrict \(\omega_f\) to the diagonal of \(X\times_YX\), move the diagonal to the zero section of \(T_{X/Y}\), and use the orientation of a vector bundle.

The projection \(p_1\) is the base change of \(f\) along \(f\), so Proposition proposition-exceptional-inverse-image-smooth-morphism shows that \(\omega_{p_1}\simeq p_2^{*}\omega_f\). Since \(p_2\circ\Delta\) is the identity of \(X\), we get the first isomorphism in

\[
\begin{aligned}
\omega_{f}
&\simeq L\Delta^{*}\omega_{p_1}
\xrightarrow{\sim}z^{*}\omega_{\pi}
\xrightarrow{\sim}z^{*}(\mathbf Z/\ell^n\mathbf Z)_{T_{X/Y}}(d)[2d]
\simeq(\mathbf Z/\ell^n\mathbf Z)_{X}(d)[2d].
\end{aligned}
\]

The second isomorphism is Lemma \hyperref[bp27]{9.9} for the section \(\Delta\) of \(p_1\), and the third is the inverse of the isomorphism of Proposition \hyperref[bp19]{9.6} at the constant sheaf, applied to the vector bundle \(T_{X/Y}\) over \(X\). The last one holds because the pullback of a constant sheaf is constant. Since each step is attached to \(f\) without choices, the composite is canonical.
\end{proof}

% text: prose-09-16
\phantomsection\label{bp36}%
The orientation (9.14) is defined without coordinates, but on a coordinate chart it agrees with the identification given by the trace map of affine space (9.11). This is what makes the map \(f^*K(d)[2d]\to f^!K\) below the expected one, and the trace map in Section \hyperref[bp72]{9.6} relies on it as well.

We prove the agreement for a vector bundle, where both orientations can be computed, and reduce a coordinate chart to the vector bundle \(\mathbf A^d_Y\) by naturality.

\emph{Vector bundles.} Let \(\pi:V\to S\) be a vector bundle of rank \(d\). Here the orientation obtained from the diagonal is the orientation of Proposition \hyperref[bp19]{9.6}, because the construction from the diagonal moves the diagonal to the zero section and then deforms the zero section to itself, and neither step changes integration along the fibres. For the first step, regard \(V\times_SV\) as a vector bundle over the first factor. The automorphism \((v,w)\mapsto(v,w-v)\) of \(V\times_SV\) over the first factor takes the diagonal to the zero section, and it preserves integration in the second factor. To see this, let the translation vector vary over \(V\). Its effect on the top compactly supported direct image, identified with the constant sheaf by integration, is then multiplication by a locally constant unit on \(V\), and this unit is \(1\) at the zero vector, where the translation is the identity. Since each geometric fibre of \(V\to S\) is connected, it is \(1\) everywhere.

Next, the deformation to the normal bundle of the zero section of \(V\) is \(V\times\mathbf A^1\), with the map \(V\times\mathbf A^1\to V\times\mathbf A^1\), \((w,t)\mapsto(tw,t)\), to the original family. This map is defined by scalar multiplication, without a trivialization, and in a local trivialization it is \(x_i=tu_i\), as in the coordinate description above. Integration on this product family is the pullback of the original integration. Therefore the comparison of Lemma \hyperref[bp27]{9.9}, which is the identity at \(t=1\), is the identity under these integrations, in particular at \(t=0\). Combined with the translation, this identifies the orientation obtained from the diagonal for \(\pi\) with the orientation of the vector bundle.

\emph{Coordinate charts.} A chart \(e:U\to\mathbf A^d_Y\) relates the diagonal of \(U\) to the diagonal of \(\mathbf A^d_Y\), where the case of vector bundles applies, so we have to transport the orientation along \(e\). By Lemma \hyperref[bp30]{9.10} it suffices to see that every ingredient of the orientation of a chart obtained from the diagonal is pulled back from affine space. So let \(U\subset X\) be open with an étale map \(e:U\to\mathbf A^d_Y\) over \(Y\). Then the étale map \(e\times e:U\times_YU\to\mathbf A^d_Y\times_Y\mathbf A^d_Y\) carries the diagonal of \(U\) into the diagonal of \(\mathbf A^d_Y\), whose inverse image is \(U\times_{\mathbf A^d_Y}U\). In it the diagonal of \(U\) is open and closed, because \(e\) is separated and étale, and removing the other components, as in the proof of Lemma \hyperref[bp27]{9.9}, we obtain a cartesian square of centers. The induced map of normal bundles is then the differential \(T_{U/Y}\simeq e^*T_{\mathbf A^d_Y/Y}\), a pullback of vector bundles. Likewise, the dualizing complex is pulled back from the affine model, because by composing adjunctions the exceptional inverse image along an étale map is its ordinary inverse image. For the comparison obtained from the deformation, we factor \(e\times e\) as

\[
U\times_YU\xrightarrow{\mathrm{id}_U\times e}
U\times_Y\mathbf A^d_Y\xrightarrow{e\times\mathrm{id}}
\mathbf A^d_Y\times_Y\mathbf A^d_Y.
\]

The second arrow is the base change of the first projection along \(e\), and it carries the diagonal section to the graph of \(e\) over \(U\), so Lemma lemma-naturality-under-base-change-etale applies to it. Part (2), over the common base \(U\), applies to the first arrow after we remove the other components of the inverse image of that graph. In this factorization the map on special fibres is the displayed differential, and the comparison at \(t=1\) is the one obtained from the adjunction in coordinates. Moreover, integration on vector bundles is preserved by the linear isomorphism of normal bundles. By the case of vector bundles, we conclude that the orientation on \(U\) obtained from the diagonal is the orientation of affine space given by its adjunction in coordinates.

Combining the orientation with the isomorphism \(f^!K\simeq f^*K\otimes^{\mathbf L}\omega_f\), we obtain the isomorphism \(\gamma_f\) of Theorem \hyperref[bp9]{9.3} with finite coefficients (Proposition \hyperref[bp37]{9.12}). We state the result for an arbitrary finite admissible ring.

% object: proposition-smooth-purity-finite-coefficients (blueprint/chap-Smooth-purity/sec-Finite-coefficients-3/proposition-smooth-purity-finite-coefficients/_section.md)
\phantomsection\label{bp37}%
\begin{bpstatementit}{Proposition~9.12}{ (Smooth purity with finite coefficients)}
Let \(A\) be any finite admissible coefficient ring, and let \(f:X\to Y\) be smooth of pure relative dimension \(d\). For \(K\in D^b_{\mathrm{ctf}}(Y_{\mathrm{pro}\acute{e}t},A)\) there is a canonical isomorphism, natural in \(K\),

\[
f_A^{*}K(d)[2d]
\xrightarrow{\ \gamma_{f,A}(K)\ }f_A^{!}K.
\tag{9.15}
\]
\end{bpstatementit}

% proof: proof/proposition-smooth-purity-finite-coefficients
\begin{proof}[{Proof}]\phantomsection\label{bp38}%
The calculation so far was made with coefficients \(\mathbf Z/\ell^n\mathbf Z\), and every ingredient in it extends along \(\mathbf Z/\ell^n\mathbf Z\to A\), where we choose \(n\) with \(\ell^nA=0\). Once the Kummer and diagonal classes are extended, the calculation works with \(A\) throughout. In particular, the splitting for the projective line has two free summands, and its integration and triangle identities keep the same normalization, since (9.9) is an identity of endomorphisms of a constant sheaf and is preserved by extension of scalars. So the local adjunctions compute \(f_A^!\), and the arguments for vector bundles and deformations then show that \(\omega_{f,A}\simeq A_X(d)[2d]\). We define \(\gamma_{f,A}(K)\) as the composite

\[
f_A^{*}K\otimes^{\mathbf L}A_X(d)[2d]
\xrightarrow{\sim}f_A^{*}K\otimes^{\mathbf L}\omega_{f,A}
\xrightarrow{\sim}f_A^{!}K.
\]

These are the isomorphisms given by the orientation and by \(\mu_{f,K}\), and both are natural in \(K\). By (9.17) below, \(\gamma_{f,A}\) does not depend on the choice of \(n\) and is compatible with maps of finite rings.
\end{proof}


% object: sec-ell-adic-coefficients-4 (blueprint/chap-Smooth-purity/sec-ell-adic-coefficients-4/_section.md), content hidden
\subsection*{\texorpdfstring{9.5\quad $\ell$-adic coefficients}{9.5 ell-adic coefficients}}
\phantomsection\label{bp39}%
\addcontentsline{toc}{subsection}{\texorpdfstring{9.5\quad $\ell$-adic coefficients}{9.5 ell-adic coefficients}}
% text: prose-09-17
\phantomsection\label{bp40}%
Let \(f:X\to Y\) be smooth of pure relative dimension \(d\). For \(\mathcal O_E\)-coefficients we want an isomorphism of functors \(f^*(-)(d)[2d]\simeq f^!\) on \(\Dbc(-,\mathcal O_E)\). Here \(f^!_{\mathcal O_E}\) is defined by Brown representability, not as a limit of the finite \(f^!\). The reduction principle of Section \bpout{7.3} does not apply, because it checks that a given morphism is an isomorphism, and we have no morphism \(f^*(-)(d)[2d]\to f^!\) with \(\mathcal O_E\)-coefficients yet. Instead we deduce Theorem \hyperref[bp9]{9.3} for \(\mathcal O_E\), \(E\) and \(\overline{\mathbf Q}_\ell\) from Proposition \hyperref[bp37]{9.12} by building the right adjoint ourselves. The key input is that the finite isomorphisms \(\gamma_{f,A}\) commute with change of finite coefficients (9.17), which lets us construct an explicit right adjoint of \(Rf_{!,\mathcal O_E}\) from towers (Proposition \hyperref[bp61]{9.19}). Inverting \(\ell\) then gives the cases \(E\) and \(\overline{\mathbf Q}_\ell\) (Proposition \hyperref[bp66]{9.20} and Theorem \hyperref[bp69]{9.21}). Each right adjoint we build is identified with \(f^!\) by the uniqueness of adjoints, normalized by the counits.

% object: proposition-sp-unique-adjoint (blueprint/chap-Smooth-purity/sec-ell-adic-coefficients-4/proposition-sp-unique-adjoint/_section.md)
\phantomsection\label{bp41}%
\begin{bpstatementit}{Proposition~9.13}{ (Uniqueness of the right adjoint of $Rf_!$)}
Let \(\Lambda\) be an admissible coefficient ring and let \(J\) be a right adjoint of \(Rf_!:D^b_{\mathrm{ctf}}(X,\Lambda)\to D^b_{\mathrm{ctf}}(Y,\Lambda)\) with counit \(\varepsilon^J\). There is a unique natural isomorphism \(\rho:f^!\xrightarrow{\sim}J\) with \(\varepsilon^J(K)\circ Rf_!(\rho(K))=\varepsilon_f(K)\).
\end{bpstatementit}

% proof: proof/proposition-sp-unique-adjoint
\begin{proof}[{Proof}]\phantomsection\label{bp42}%
Let \(\rho(K)\) be the \(J\)-transpose of \(\varepsilon_f(K)\) and \(\rho'(K)\) the \(f^!\)-transpose of \(\varepsilon^J(K)\). Since \(\rho'\rho\) and the identity have the same transpose, they are equal, and by symmetry \(\rho\rho'=\mathrm{id}\). For uniqueness, note that a map compatible with the counits has \(J\)-transpose \(\varepsilon_f(K)\), so it is \(\rho(K)\).
\end{proof}

% object: subsec-Change-of-finite-coefficients (blueprint/chap-Smooth-purity/sec-ell-adic-coefficients-4/subsec-Change-of-finite-coefficients/_section.md), content hidden
\subsubsection*{\texorpdfstring{9.5.1\quad Change of finite coefficients}{9.5.1 Change of finite coefficients}}
\phantomsection\label{bp43}%
\addcontentsline{toc}{subsubsection}{\texorpdfstring{9.5.1\quad Change of finite coefficients}{9.5.1 Change of finite coefficients}}
% text: prose-09-18
\phantomsection\label{bp44}%
The levels \(\mathcal O_E/\ell^n\) of a tower are related by extension of scalars. The first thing to know is therefore how \(Rf_!\) and \(f^!\) behave under a map of finite rings \(A\to B\). The comparison for \(f^!\) is the transpose of a counit, and the counits commute with \(-\otimes^{\mathbf L}B\), because on a chart they are built from the trace of the projective line and sums over fibres.

Let \(A\to B\) be a map of finite commutative rings killed by a power of \(\ell\).

% object: definition-sp-scalar-extension (blueprint/chap-Smooth-purity/sec-ell-adic-coefficients-4/subsec-Change-of-finite-coefficients/definition-sp-scalar-extension/_section.md)
\phantomsection\label{bp45}%
\begin{bpstatement}{Definition~9.14}{ (Extension of scalars)}
For a scheme \(T\) and \(K\in D^b_{\mathrm{ctf}}(T_{\mathrm{pro}\acute{e}t},A)\) put \(K_B=K\otimes^{\mathbf L}_{A_T}B_T\in D^b_{\mathrm{ctf}}(T_{\mathrm{pro}\acute{e}t},B)\), and for a morphism \(u\) put \(u_B=u\otimes^{\mathbf L}\mathrm{id}_{B_T}\). For a further map \(B\to C\) we have \((K_B)_C\simeq K_C\) by associativity of the tensor product.
\end{bpstatement}

% text: prose-09-19
\phantomsection\label{bp46}%
With this notation, the comparison for \(f^!\) is characterized by its counit: it is the unique map that makes the counit over \(B\) agree with the extended counit over \(A\).

% object: lemma-sp-reduction (blueprint/chap-Smooth-purity/sec-ell-adic-coefficients-4/subsec-Change-of-finite-coefficients/lemma-sp-reduction/_section.md)
\phantomsection\label{bp47}%
\begin{bpstatementit}{Lemma~9.15}{ (Extension of scalars commutes with $Rf_!$ and $f^!$)}
For \(M\in D^b_{\mathrm{ctf}}(X_{\mathrm{pro}\acute{e}t},A)\) and \(K\in D^b_{\mathrm{ctf}}(Y_{\mathrm{pro}\acute{e}t},A)\) there are natural isomorphisms \(b(M):(Rf_{!,A}M)_B\xrightarrow{\sim}Rf_{!,B}(M_B)\) and \(c_{A,B}(K):(f_A^!K)_B\xrightarrow{\sim}f_B^!(K_B)\), where \(c_{A,B}(K)\) is determined by

\[
\varepsilon_B(K_B)\circ Rf_{!,B}\bigl(c_{A,B}(K)\bigr)
=(\varepsilon_A(K))_B\circ b(f_A^!K)^{-1}.
\tag{9.16}
\]

For \(A\to B\to C\), the comparisons for \(A\to C\) are the composites of those for \(A\to B\) and \(B\to C\).
\end{bpstatementit}

% proof: proof/lemma-sp-reduction
\begin{proof}[{Proof}]\phantomsection\label{bp48}%
The comparison \(b\) for \(Rf_!\) is a direct computation, and \(c_{A,B}\) is then determined by the adjunction. To see that \(c_{A,B}\) is invertible we use the explicit form of \(f^!\) on a chart, where it is an ordinary inverse image.

The map \(b\) is the change of coefficients for \(Rf_!\) (Proposition \bpout{7.10}). To see that it is an isomorphism, note that extension by zero commutes with \(-\otimes^{\mathbf L}_AB\), and for \(R\bar f_*\) we resolve \(B\) by finite free \(A\)-modules and use the flat models of Lemma \bpout{3.13}. The comparison holds on each free term, which is a finite direct sum of copies of \(A\), and since \(R\bar f_*\) has finite cohomological dimension, only finitely many terms of the resolution contribute in each degree, so the comparison holds for every \(B\). Next, we define \(c_{A,B}(K)\) as the transpose of the right-hand side of (9.16). Whether it is an isomorphism is a local question, so we may restrict to a chart \(f|_U=p_d\circ e\) of Proposition proposition-exceptional-inverse-image-smooth-morphism. There the counit is built from the trace of the projective line and sums over the fibres of \(e\), which commute with extension of coefficients (Proposition \bpout{6.7}). On the chart, therefore, \(c_{A,B}(K)\) is the transpose of the comparison for pullback \((e^*p_d^*K(d)[2d])_B\simeq e^*p_d^*(K_B)(d)[2d]\), which is an isomorphism. Hence \(c_{A,B}(K)\) is an isomorphism on every chart. Transitivity holds because both sides have the same transpose: for \(A\to B\to C\), the counit identity (9.16) for \(A\to C\) is obtained by applying the one for \(B\to C\) after the one for \(A\to B\).
\end{proof}

% text: prose-09-20
\phantomsection\label{bp49}%
It follows that the isomorphisms \(\gamma_{f,A}\) are compatible with change of coefficients: for \(K\in D^b_{\mathrm{ctf}}(Y_{\mathrm{pro}\acute{e}t},A)\),

\[
c_{A,B}(K)\circ(\gamma_{f,A}(K))_B=\gamma_{f,B}(K_B),
\tag{9.17}
\]

the sources being identified by \((f_A^*K(d)[2d])_B\simeq f_B^*(K_B)(d)[2d]\). Indeed, \(\gamma_{f,A}\) is the map \(\mu_{f,K}\) composed with the inverse of the orientation, and each factor commutes with change of coefficients.

The map \(\mu_{f,K}\) is built from the projection formula and the counit. The projection formula commutes with extension of scalars, which is itself a tensor product, and (9.16) relates the counits over \(A\) and \(B\). Hence \(\mu_{f,K}\) commutes with change of coefficients.

The orientation is a map between invertible sheaves placed in degree \(-2d\), so its compatibility is local. For vector bundles it follows from the construction via the projective line (Proposition \hyperref[bp19]{9.6}), since the Kummer class and the splitting (9.8) are compatible with extension of scalars. On the deformation family we use the uniqueness in Lemma \hyperref[bp27]{9.9} over the base of the family. It applies once we know that the extended comparison is a comparison of the kind considered there, and this holds because \(c\) commutes with smooth exceptional base change and with composition with étale maps. Indeed, both sides are transposes of counits, and by the proof of Lemma \hyperref[bp47]{9.15} they agree on charts. Since the original comparison is the identity at \(t=1\), so is the extended one, and by uniqueness the extended comparison is the comparison of Lemma \hyperref[bp27]{9.9} over \(B\).


% object: subsec-Integral-coefficients-2 (blueprint/chap-Smooth-purity/sec-ell-adic-coefficients-4/subsec-Integral-coefficients-2/_section.md), content hidden
\subsubsection*{\texorpdfstring{9.5.2\quad Integral coefficients}{9.5.2 Integral coefficients}}
\phantomsection\label{bp50}%
\addcontentsline{toc}{subsubsection}{\texorpdfstring{9.5.2\quad Integral coefficients}{9.5.2 Integral coefficients}}
% text: prose-09-21
\phantomsection\label{bp51}%
Fix a finite extension \(E/\mathbf Q_\ell\). An object \(M\) of \(\Dbc(T,\mathcal O_E)\) is derived complete, \(M\simeq R\varprojlim_nM/\ell^n\), and the finite theory applies to each \(M/\ell^n\). The units, counits and transition maps at the finite levels, however, live in the different categories \(D(T_{\mathrm{pro}\acute{e}t},\mathcal O_E/\ell^n)\), and a family of them that commutes with the transition maps only up to unspecified isomorphisms does not define a morphism between the derived limits. We therefore work in the derived category \(D(\operatorname{Tow}(T))\) of towers \((C_n)_n\), where \(C_n\) is a sheaf of \(\mathcal O_{E,T}/\ell^n\)-modules on \(T_{\mathrm{pro}\acute{e}t}\) and the transition maps are linear over \(\mathcal O_{E,T}/\ell^{n+1}\to\mathcal O_{E,T}/\ell^n\). These are the modules over \((\mathcal O_{E,T}/\ell^n)_n\) in the topos of inverse systems [Stacks, \href{https://stacks.math.columbia.edu/tag/0940}{0940}]. In this category a compatible family of maps is one morphism. Evaluation \(C_\bullet\mapsto C_n\) is exact, and the evaluations jointly detect isomorphisms, but they do not detect equality of morphisms. This is why the maps below are always compared in \(D(\operatorname{Tow}(T))\). A tower is \emph{constructible} if each \(C_n\) lies in \(D^b_{\mathrm{ctf}}(T_{\mathrm{pro}\acute{e}t},\mathcal O_E/\ell^n)\) with bounds independent of \(n\), and \emph{normalized} if its reduction maps \(C_n\otimes^{\mathbf L}\mathcal O_{E,T}/\ell^m\to C_m\) are isomorphisms for all \(m\le n\).

We need two facts about integral complexes. Part (1) of the next proposition gives the uniform bounds in the definition of a constructible tower, and part (2) shows that the derived limit of such a tower is again an object of \(\Dbc\).

% object: proposition-sp-integral-towers (blueprint/chap-Smooth-purity/sec-ell-adic-coefficients-4/subsec-Integral-coefficients-2/proposition-sp-integral-towers/_section.md)
\phantomsection\label{bp52}%
\begin{bpstatementit}{Proposition~9.16}{ (Integral constructible complexes as towers)}
\begin{enumerate}
\item For \(K\in\Dbc(T,\mathcal O_E)\) there are integers \(a\le b\) such that every reduction \(K/\ell^n\) lies in \(D^b_{\mathrm{ctf}}(T_{\mathrm{pro}\acute{e}t},\mathcal O_E/\ell^n)\), with cohomology and Tor amplitude in \([a,b]\).
\item If \(K\in D(T_{\mathrm{pro}\acute{e}t},\mathcal O_E)\) is derived complete and \(K/\ell\) lies in \(D^b_{\mathrm{ctf}}(T_{\mathrm{pro}\acute{e}t},\mathcal O_E/\ell)\), then \(K\in\Dbc(T,\mathcal O_E)\).
\end{enumerate}
\end{bpstatementit}

% proof: proof/proposition-sp-integral-towers
\begin{proof}[{Proof}]\phantomsection\label{bp53}%
For (1), by Lemma \bpout{5.19}, \(K\) is a finite extension of shifted sheaves \(a_!C\) with \(C\) of Tor amplitude \([-1,0]\) on a stratum. For a single such piece the reductions \((a_!C)/\ell^n=a_!(C/\ell^n)\) are constructible with Tor amplitude in \([-1,0]\) for every \(n\), since reduction commutes with extension by zero. Moreover, Tor amplitude and constructibility of the reductions are preserved by extensions and shifts, and since the number of steps and the shifts do not depend on \(n\), the bounds do not either. For (2), we reduce further from \(\ell\) to the uniformizer. The reduction \(K/\pi_E=(K/\ell)\otimes^{\mathbf L}_{\mathcal O_E/\ell}\mathcal O_E/\pi_E\) is bounded constructible, because \(K/\ell\) has finite Tor amplitude over \(\mathcal O_E/\ell\). Since \(K\) is derived complete, \(K\in\Dbc(T,\mathcal O_E)\) by the criterion (2)\(\Rightarrow\)(1) of Proposition \bpout{3.11}.
\end{proof}

% text: prose-09-22
\phantomsection\label{bp54}%
Nothing is lost by passing from \(\mathcal O_E\)-complexes to towers: a constructible \(\mathcal O_E\)-complex is recovered from its tower of reductions as the derived limit, and the towers that arise in this way are exactly the normalized constructible ones.

% object: lemma-sp-tower-adjunction-foundation (blueprint/chap-Smooth-purity/sec-ell-adic-coefficients-4/subsec-Integral-coefficients-2/lemma-sp-tower-adjunction-foundation/_section.md)
\phantomsection\label{bp55}%
\begin{bpstatementit}{Lemma~9.17}{ (Reduction to towers is left adjoint to the derived limit)}
The functor \(\operatorname{red}_TM=(M/\ell^n)_n\) is left adjoint to \(R\varprojlim:D(\operatorname{Tow}(T))\to D(T_{\mathrm{pro}\acute{e}t},\mathcal O_E)\), with unit the derived completion map \(M\to R\varprojlim_nM/\ell^n\). This adjunction restricts to an equivalence between \(\Dbc(T,\mathcal O_E)\) and the normalized constructible towers.
\end{bpstatementit}

% proof: proof/lemma-sp-tower-adjunction-foundation
\begin{proof}[{Proof}]\phantomsection\label{bp56}%
The adjunction is a composite of two standard adjunctions of ringed topoi. For the equivalence we check that the unit is an isomorphism on \(\Dbc\) and that the counit is an isomorphism on normalized constructible towers. The second is the substantial point, since there we must show that the limit of such a tower is again constructible.

The adjunction is \(Lh^*p^{-1}\dashv Rp_*h_*\), for the change of rings \(h\) from towers to constant inverse systems of \(\mathcal O_{E,T}\)-modules and the morphism \(p\) to \(T_{\mathrm{pro}\acute{e}t}\) with \(p_*=\varprojlim\) [Stacks, \href{https://stacks.math.columbia.edu/tag/06YV}{06YV}, \href{https://stacks.math.columbia.edu/tag/07A6}{07A6}]. By Proposition proposition-sp-integral-towers, \(\operatorname{red}_T\) maps \(\Dbc(T,\mathcal O_E)\) into normalized constructible towers. On \(\Dbc\) the unit is an isomorphism, since objects of \(\Dbc\) are derived complete.

Conversely, let \(C_\bullet\) be normalized and put \(C=R\varprojlim C_n\). The key point is that \(C/\ell^m\simeq C_m\). Indeed, the complex \(C/\ell^m\) differs from \(C_m\) by the derived limit of the tower \((C_m[1])_n\) with transition maps multiplication by \(\ell\), and this limit vanishes, because \(\ell^m\) kills \(C_m\) and so any \(m\) consecutive transition maps compose to zero. Consequently \(C\) is derived complete and \(C/\ell\simeq C_1\) is constructible. Moreover \(C\) is bounded, because the \(C_n\) have cohomology in a fixed interval, and on the replete topos \(T_{\mathrm{pro}\acute{e}t}\) derived limits of such towers increase the range by at most one [Stacks, \href{https://stacks.math.columbia.edu/tag/0940}{0940}]. Therefore \(C\in\Dbc(T,\mathcal O_E)\) by Proposition proposition-sp-integral-towers. Finally, the counit \(\operatorname{red}_TC\to C_\bullet\) is at level \(m\) the reduction map \(C/\ell^m\to C_m\), which is an isomorphism by what we just showed. Since evaluations detect isomorphisms, the counit is an isomorphism.
\end{proof}

% text: prose-09-23
\phantomsection\label{bp57}%
All operations on towers are defined levelwise, and the functor \(Rg_*\) is computed levelwise as well, since an injective tower is levelwise a summand of finite products of injective sheaves. So proper base change and the projection formula hold for towers.

Moreover \(Rf_{!,\bullet}\) preserves constructible towers. Each level lies in \(D^b_{\mathrm{ctf}}\) by Proposition \bpout{7.1}, so we only check that the bounds are independent of \(n\). For cohomology, note that for a compactification \(f=pj\), \(Rf_!\) has amplitude \([0,\operatorname{cd}_\ell(p)]\) for every torsion coefficient ring [Stacks, \href{https://stacks.math.columbia.edu/tag/0F10}{0F10}]. For Tor amplitude we use the projection formula \((Rf_!C_n)\otimes^{\mathbf L}N\simeq Rf_!(C_n\otimes^{\mathbf L}N)\) for every \(\mathcal O_E/\ell^n\)-module \(N\), which follows from the comparison \(b\) and shows that if all \(C_n\) have Tor amplitude in \([a,b]\), then all \(Rf_!C_n\) have Tor amplitude in \([a,b+\operatorname{cd}_\ell(p)]\).

We do not use Brown representability for towers. Instead we construct the right adjoint of \(Rf_{!,\bullet}\) on constructible towers by repeating the geometric argument with finite coefficients in the category of towers.

% object: proposition-sp-tower-geometric (blueprint/chap-Smooth-purity/sec-ell-adic-coefficients-4/subsec-Integral-coefficients-2/proposition-sp-tower-geometric/_section.md)
\phantomsection\label{bp58}%
\begin{bpstatementit}{Proposition~9.18}{ (The adjunction $Rf_!\dashv f^!$ on constructible towers)}
There is an adjunction

\[
Rf_{!,\bullet}\dashv f_\bullet^!
\quad\text{between the constructible towers on }X\text{ and on }Y.
\tag{9.18}
\]

At level \(n\) its unit and counit are those of the finite adjunction, so \((f_\bullet^!C_\bullet)_n\simeq f^!_{\mathcal O_E/\ell^n}C_n\), and the reduction maps of \(f_\bullet^!C_\bullet\) are \(f^!\) of those of \(C_\bullet\) composed with the comparisons \(c\) of Lemma \hyperref[bp47]{9.15}.
\end{bpstatementit}

% proof: proof/proposition-sp-tower-geometric
\begin{proof}[{Proof}]\phantomsection\label{bp59}%
We repeat the construction of Section \hyperref[bp15]{9.3} in the topos of towers. The construction has two parts: on a chart we write down an explicit adjunction, built from the trace of the projective line, and then we glue these local adjunctions over a finite cover. The gluing is done with pairs consisting of an object and a counit, as in Definition \hyperref[bp17]{9.5}, because such pairs are unique up to unique isomorphism and so need no choice of transition maps.

The Kummer sequence \(0\to(\mu_{\ell^n})_n\to(\mathbf G_m,(\cdot)^\ell)\to(\mathbf G_m,\mathrm{id})\to0\) is exact in towers. Hence Chern classes, the diagonal class and the trace of the projective line are single morphisms of towers. Moreover, the identities (9.9) are equalities of maps between sheaves of towers in degree zero, so they can be checked at each level. The construction of (9.10) then gives the adjunction for the projective line in \(D(\operatorname{Tow})\), and extension by zero and the étale adjunction give local adjunctions on charts.

We glue the local adjunctions in \(D(\operatorname{Tow})\), by induction on a finite cover. Let \(X=U\cup V\), and let \((H_U,\varepsilon_U)\) and \((H_V,\varepsilon_V)\) be pairs representing \(\operatorname{Hom}(Rf_{!,\bullet}(-),K_\bullet)\) on \(U\) and on \(V\). Put \(W=U\cap V\). Both pairs restrict to pairs representing this functor on \(W\), so by Yoneda there is a unique counit-preserving isomorphism between the restrictions, and we call the common restriction \(H_W\). Let \(H\) be the cone of \(j_{W!}H_W\to j_{U!}H_U\oplus j_{V!}H_V\), as in the triangle (9.12). Since the counits \(\varepsilon_U,\varepsilon_V\) agree on \(Rf_{!,\bullet}j_{W!}H_W\), the exact Hom sequence shows that they lift to a counit \(\varepsilon:Rf_{!,\bullet}H\to K_\bullet\). The five lemma, applied as in the finite case, then shows that \((H,\varepsilon)\) represents \(\operatorname{Hom}(Rf_{!,\bullet}(-),K_\bullet)\) on \(X\). The tower \(H\) is also constructible. Indeed, \(H|_U\simeq H_U\) and \(H|_V\simeq H_V\), so the union of their intervals bounds its cohomology and Tor amplitude uniformly in \(n\).

Repeating this step over a finite cover of \(X\) by charts gives a pair representing \(\operatorname{Hom}(Rf_{!,\bullet}(-),K_\bullet)\) on all of \(X\), and we let \(f^!_\bullet K_\bullet\) be its object. By uniqueness of representing pairs, this is functorial in \(K_\bullet\) and gives the adjunction (9.18).

The levels are the finite ones, because evaluation at level \(n\) carries the glued pair to a pair representing \(\operatorname{Hom}(Rf_!(-),K_n)\) with coefficients \(\mathcal O_E/\ell^n\). Indeed, on each chart it is the local finite adjunction, and the five lemma argument applies at level \(n\), so by uniqueness this pair is the finite one. As the counit is one morphism of towers, it commutes with the transition maps, and transposing this identity gives the reduction maps.
\end{proof}

% text: prose-09-24
\phantomsection\label{bp60}%
The orientation passes to towers as well. By the finite case each level of \(\omega_{f,\bullet}[-2d]\) is an invertible sheaf, so it is an invertible sheaf of towers. Moreover, the comparisons for vector bundles and for the deformation family are maps of sheaves of towers in degree zero, and equality of such maps can be tested level by level, unlike equality of maps in \(D(\operatorname{Tow})\) in general. Hence the uniqueness in Lemma \hyperref[bp27]{9.9} holds for them, because it holds at each level. Composing the inverse of the orientation with the map \(\mu_{f,K}\) on towers, we obtain an isomorphism, natural in the constructible tower \(K_\bullet\),

\[
\gamma_{f,\bullet}(K_\bullet):f_\bullet^*K_\bullet(d)[2d]\xrightarrow{\sim}f_\bullet^!K_\bullet,
\tag{9.19}
\]

whose level \(n\) is \(\gamma_{f,\mathcal O_E/\ell^n}(K_n)\).

Both functors in (9.18) preserve normalized towers. For \(Rf_{!,\bullet}\) this follows from the comparison \(b\), and for \(f_\bullet^!\) from (9.17) and (9.19).

Since \(R\bar f_*\) and \(j_!\) commute with \(R\varprojlim\) [Stacks, \href{https://stacks.math.columbia.edu/tag/09BP}{09BP}], we have

\[
Rf_{!,\mathcal O_E}M\simeq R\varprojlim_nRf_{!,\mathcal O_E/\ell^n}(M/\ell^n)
\]

for \(M\in\Dbc(X,\mathcal O_E)\). For \(K\in\Dbc(Y,\mathcal O_E)\) we put

\[
H_{f,\mathcal O_E}(K)=R\varprojlim_n f_{\mathcal O_E/\ell^n}^!(K/\ell^n)\in\Dbc(X,\mathcal O_E),
\tag{9.20}
\]

the derived limit of the normalized tower \(f_\bullet^!\operatorname{red}_YK\).

The functor \(H_{f,\mathcal O_E}\) is given by an explicit formula, whereas we only know that \(f^!_{\mathcal O_E}\) exists. The next proposition shows that the formula defines a right adjoint, and so, by uniqueness, computes \(f^!_{\mathcal O_E}\).

% object: proposition-sp-integral-adjunction (blueprint/chap-Smooth-purity/sec-ell-adic-coefficients-4/subsec-Integral-coefficients-2/proposition-sp-integral-adjunction/_section.md)
\phantomsection\label{bp61}%
\begin{bpstatementit}{Proposition~9.19}{ (The integral right adjoint built from towers)}
With the unit and counit obtained from towers, \(H_{f,\mathcal O_E}\) is right adjoint to \(Rf_{!,\mathcal O_E}\) on \(\Dbc(-,\mathcal O_E)\). Hence there is a unique counit-preserving isomorphism \(H_{f,\mathcal O_E}\simeq f^!_{\mathcal O_E}\) with the right adjoint of Theorem \bpout{5.31}.
\end{bpstatementit}

% proof: proof/proposition-sp-integral-adjunction
\begin{proof}[{Proof}]\phantomsection\label{bp62}%
We move the adjunction on towers down to \(\mathcal O_E\)-complexes along the equivalence of Lemma \hyperref[bp55]{9.17}. For this we need both functors to stay inside normalized towers, which is where the equivalence holds.

Since \(Rf_{!,\bullet}\) and \(f^!_\bullet\) preserve normalized towers, the equivalence of Lemma \hyperref[bp55]{9.17} gives \(\operatorname{red}_Y(Rf_{!,\mathcal O_E}M)\simeq Rf_{!,\bullet}\operatorname{red}_XM\) and \(\operatorname{red}_XH_{f,\mathcal O_E}(K)\simeq f^!_\bullet\operatorname{red}_YK\). Using this lemma and (9.18), we get natural bijections

\[
\begin{aligned}
\Hom(Rf_{!,\mathcal O_E}M,K)
&\simeq\Hom(Rf_{!,\bullet}\operatorname{red}_XM,\operatorname{red}_YK)\\
&\simeq\Hom(\operatorname{red}_XM,f_\bullet^!\operatorname{red}_YK)
\simeq\Hom(M,H_{f,\mathcal O_E}(K)).
\end{aligned}
\]

The second statement follows from Proposition \hyperref[bp41]{9.13}.
\end{proof}

% text: prose-09-25
\phantomsection\label{bp63}%
Applying \(R\varprojlim\) to (9.19) at \(\operatorname{red}_YK\), we obtain \(\gamma_{f,\mathcal O_E}:f^*_{\mathcal O_E}K(d)[2d]\xrightarrow{\sim}f^!_{\mathcal O_E}K\), whose reductions are the isomorphisms \(\gamma_{f,\mathcal O_E/\ell^n}\). This gives the isomorphism (9.2) for \(\mathcal O_E\), and for \(E=\mathbf Q_\ell\) the one for \(\mathbf Z_\ell\).


% object: subsec-Field-coefficients (blueprint/chap-Smooth-purity/sec-ell-adic-coefficients-4/subsec-Field-coefficients/_section.md), content hidden
\subsubsection*{\texorpdfstring{9.5.3\quad Field coefficients}{9.5.3 Field coefficients}}
\phantomsection\label{bp64}%
\addcontentsline{toc}{subsubsection}{\texorpdfstring{9.5.3\quad Field coefficients}{9.5.3 Field coefficients}}
% text: prose-09-26
\phantomsection\label{bp65}%
With coefficients in \(E\) there is no tower to pass to the limit over, but none is needed. Indeed, an object of \(\Dbc(T,E)\) comes from an integral complex by inverting \(\ell\), and its Hom groups are the integral Hom groups tensored with \(E\) (9.21). Since an adjunction over \(\mathcal O_E\) is a natural bijection of Hom groups, tensoring it with \(E\) gives an adjunction over \(E\) (Proposition \hyperref[bp66]{9.20}). For \(\overline{\mathbf Q}_\ell\) the Hom groups are colimits over finite extensions of \(\mathbf Q_\ell\), and the adjunctions over these fields pass to the colimit (Theorem \hyperref[bp69]{9.21}). For each coefficient ring we thus construct an explicit right adjoint of \(Rf_!\), and we identify it with \(f^!\) by the uniqueness of adjoints (Proposition \hyperref[bp41]{9.13}).

Every object of \(\Dbc(T,E)\) is \(K_0[1/\ell]\) for some \(K_0\in\Dbc(T,\mathcal O_E)\), called an integral model, and for integral models

\[
\Hom(K_0,L_0[q])\otimes_{\mathcal O_E}E
\xrightarrow{\sim}\Hom_{\Dbc(T,E)}\bigl(K_0[1/\ell],L_0[1/\ell][q]\bigr).
\tag{9.21}
\]

Both statements are Proposition \bpout{3.27}: its essential surjectivity provides the models, and (9.21) is its full faithfulness. To see the latter directly, note that the target is \(\Hom(K_0,L_0[1/\ell][q])\) by adjunction, and \(L_0[1/\ell]\) is the colimit of \(L_0\xrightarrow{\ell}L_0\xrightarrow{\ell}\cdots\). The left side is the colimit of the same system of Hom groups, so the claim is that Hom out of \(K_0\) commutes with this colimit. Since both sides are cohomological functors of \(K_0\), by the five lemma and dévissage we may take \(K_0=a_!C\) with \(C\) dualizable on a stratum (Lemma lemma-structure-integral-constructible-complexes). The stratum formula of Lemma lemma-sections-support-closed-subscheme, whose proof uses only that \(C\) is dualizable, turns Hom out of \(a_!C\) into cohomology on a quasi-compact quasi-separated scheme, and this commutes with the colimit by Corollary \bpout{2.29}. On models \(f^*\) and \(Rf_!\) are the integral functors with \(\ell\) inverted, again because \(R\bar f_*\) commutes with the colimit along \(\ell\).

With (9.21) the adjunction for \(E\) becomes a matter of inverting \(\ell\) in the integral adjunction.

% object: proposition-sp-rational (blueprint/chap-Smooth-purity/sec-ell-adic-coefficients-4/subsec-Field-coefficients/proposition-sp-rational/_section.md)
\phantomsection\label{bp66}%
\begin{bpstatementit}{Proposition~9.20}{ (Smooth purity with coefficients in $E$)}
Let \(E/\mathbf Q_\ell\) be finite. Then \(Rf_{!,E}\dashv\bigl(K\mapsto f_E^*K(d)[2d]\bigr)\) as functors between \(\Dbc(X,E)\) and \(\Dbc(Y,E)\).
\end{bpstatementit}

% proof: proof/proposition-sp-rational
\begin{proof}[{Proof}]\phantomsection\label{bp67}%
An adjunction is a natural bijection of Hom groups, and by (9.21) every Hom group over \(E\) is an integral Hom group tensored with \(E\). So we compute both sides on integral models. Choose integral models \(M_0\) and \(K_0\), so that \(M=M_0[1/\ell]\) and \(K=K_0[1/\ell]\). By (9.21) twice and Proposition \hyperref[bp61]{9.19} with \(\gamma_{f,\mathcal O_E}\),

\[
\begin{aligned}
\Hom(Rf_{!,E}M,K)&\simeq\Hom(Rf_{!,\mathcal O_E}M_0,K_0)\otimes E\\
&\simeq\Hom(M_0,f^*_{\mathcal O_E}K_0(d)[2d])\otimes E
\simeq\Hom(M,f_E^*K(d)[2d]).
\end{aligned}
\]

These bijections are given by the integral unit and counit with \(\ell\) inverted, so they are natural and do not depend on the models.
\end{proof}

% text: prose-09-27
\phantomsection\label{bp68}%
Since \(f^*(-)(d)[2d]\) is now a right adjoint of \(Rf_{!,E}\), as is \(f^!_E\), Proposition \hyperref[bp41]{9.13} gives an isomorphism \(f^!_E\xrightarrow{\sim}f_E^*(-)(d)[2d]\) compatible with the counits, and we let \(\gamma_{f,E}:f_E^*(-)(d)[2d]\xrightarrow{\sim}f^!_E\) be its inverse.

Hom groups over \(\overline{\mathbf Q}_\ell\) are colimits of Hom groups over finite extensions, so the adjunctions over the finite extensions pass to the colimit provided they correspond under the transition maps. To see that they do, let \(E\subset F\) be a finite extension. Then extension of scalars along the finite free map \(\mathcal O_E\to\mathcal O_F\) commutes with \(R\varprojlim\), and at level \(n\) it is extension along \(\mathcal O_E/\ell^n\to\mathcal O_F/\ell^n\), so we may compare the counits in \(D(\operatorname{Tow})\). Extending scalars in the pair of towers \((f_\bullet^!K_\bullet,\varepsilon)\) for \(\mathcal O_E\) gives a pair over \(\mathcal O_F\), whose levels represent \(\operatorname{Hom}(Rf_!(-),K_n)\) by Lemma \hyperref[bp47]{9.15}. By the five lemma, applied levelwise, the pair then represents \(\operatorname{Hom}(Rf_{!,\bullet}(-),K_\bullet)\), and by uniqueness it is isomorphic to the corresponding pair for \(\mathcal O_F\) by a unique counit-preserving isomorphism. Applying \(R\varprojlim\) and inverting \(\ell\), we see that the adjunctions over \(E\) and over \(F\) correspond under extension of scalars, transitively in \(E\subset F\subset F'\).

% object: theorem-sp-qbar (blueprint/chap-Smooth-purity/sec-ell-adic-coefficients-4/subsec-Field-coefficients/theorem-sp-qbar/_section.md)
\phantomsection\label{bp69}%
\begin{bpstatementit}{Theorem~9.21}{ (Smooth purity with coefficients in $\overline{\mathbf Q}_\ell$)}
We have \(Rf_{!,\overline{\mathbf Q}_\ell}\dashv\bigl(K\mapsto f_{\overline{\mathbf Q}_\ell}^*K(d)[2d]\bigr)\) as functors between \(\Dbc(X,\overline{\mathbf Q}_\ell)\) and \(\Dbc(Y,\overline{\mathbf Q}_\ell)\).
\end{bpstatementit}

% proof: proof/theorem-sp-qbar
\begin{proof}[{Proof}]\phantomsection\label{bp70}%
As for \(E\), we compare Hom groups, now as colimits over finite extensions instead of after inverting \(\ell\). By Proposition \bpout{7.11}, \(\Dbc(T,\overline{\mathbf Q}_\ell)\) is the \(2\)-colimit of the categories \(\Dbc(T,F)\) over finite extensions \(F\) of \(\mathbf Q_\ell\), so its Hom groups are the colimits of those over \(F\). The functors \(Rf_!\) and \(f^*\) commute with extension of scalars (Section \bpout{7.2}). Choosing models over some \(E\), the adjunctions of Proposition \hyperref[bp66]{9.20} over the fields \(F\supset E\) are compatible with the transition maps, and their colimit is the asserted bijection.
\end{proof}

% text: prose-09-28
\phantomsection\label{bp71}%
Applying Proposition \hyperref[bp41]{9.13} once more and inverting the resulting isomorphism, we get \(\gamma_{f,\overline{\mathbf Q}_\ell}:f_{\overline{\mathbf Q}_\ell}^*(-)(d)[2d]\xrightarrow{\sim}f^!_{\overline{\mathbf Q}_\ell}\), whose counit is obtained by extension of scalars from a finite extension of \(\mathbf Q_\ell\). What is left of Theorem \hyperref[bp9]{9.3} are the compatibilities with change of coefficients, base change and composition, and the trace normalization.



% object: sec-Compatibilities (blueprint/chap-Smooth-purity/sec-Compatibilities/_section.md), content hidden
\subsection*{\texorpdfstring{9.6\quad Compatibilities}{9.6 Compatibilities}}
\phantomsection\label{bp72}%
\addcontentsline{toc}{subsection}{\texorpdfstring{9.6\quad Compatibilities}{9.6 Compatibilities}}
% text: prose-09-29
\phantomsection\label{bp73}%
The isomorphism \(\gamma_f\) now exists for every admissible coefficient ring, and it remains to prove the properties (2)–(4) of Theorem \hyperref[bp9]{9.3}. In topology, orientations of a submersion restrict to fibres and multiply under composition, and Propositions \hyperref[bp76]{9.22} and \hyperref[bp79]{Proposition 9.23} are the corresponding statements. These propositions also let us define the trace map \(\operatorname{Tr}_f\) (Definition \hyperref[bp83]{9.24}), the analogue of integration over the fibres, on which Poincaré duality rests. With finite coefficients the compatibilities follow from the construction, since every ingredient of the orientation commutes with base change. Passing them to \(\mathbf Z_\ell\), however, is harder, because there the base change map itself has to be built as a limit.

% object: subsec-Base-change-and-composition (blueprint/chap-Smooth-purity/sec-Compatibilities/subsec-Base-change-and-composition/_section.md), content hidden
\subsubsection*{\texorpdfstring{9.6.1\quad Base change and composition}{9.6.1 Base change and composition}}
\phantomsection\label{bp74}%
\addcontentsline{toc}{subsubsection}{\texorpdfstring{9.6.1\quad Base change and composition}{9.6.1 Base change and composition}}
% text: prose-09-30
\phantomsection\label{bp75}%
Fix a cartesian square

\[
\begin{CD}
X' @>{g'}>> X\\
@V{f'}VV @VV{f}V\\
Y' @>{g}>> Y,
\end{CD}
\]

with \(f\) smooth of pure relative dimension \(d\), and a map \(A\to B\) of finite commutative rings killed by a power of \(\ell\). A subscript \(B\) denotes extension of scalars, as in Definition \hyperref[bp45]{9.14}.

\textbf{Base change.} The statement of Proposition \hyperref[bp76]{9.22} needs an exceptional base change isomorphism for \(\mathbf Z_\ell\). It is obtained as a limit of the finite ones, and the limit exists once the finite isomorphisms commute with the reduction maps, that is, once base change commutes with change of coefficients. To state this we name the four comparison maps involved. For \(M\in D^b_{\mathrm{ctf}}(X_{\mathrm{pro}\acute{e}t},A)\) and \(K\in D^b_{\mathrm{ctf}}(Y_{\mathrm{pro}\acute{e}t},A)\), let \(b^f(M):(Rf_{!,A}M)_B\xrightarrow{\sim}Rf_{!,B}M_B\) and \(c^f(K):(f_A^!K)_B\xrightarrow{\sim}f_B^!K_B\) be the change-of-coefficients maps of Lemma \hyperref[bp47]{9.15}. With \(g'^*:D^b_{\mathrm{ctf}}(X_{\mathrm{pro}\acute{e}t},A)\to D^b_{\mathrm{ctf}}(X'_{\mathrm{pro}\acute{e}t},A)\) and \(g^*:D^b_{\mathrm{ctf}}(Y_{\mathrm{pro}\acute{e}t},A)\to D^b_{\mathrm{ctf}}(Y'_{\mathrm{pro}\acute{e}t},A)\), let \(\beta_{f,A}(M):g^*Rf_{!,A}M\xrightarrow{\sim}Rf'_{!,A}g'^*M\) be the base change map for compact support, and let \(e_{f,A}(K):g'^*f_A^!K\xrightarrow{\sim}f'^!_Ag^*K\) be the exceptional base change isomorphism of Proposition proposition-exceptional-inverse-image-smooth-morphism. The claim is

\[
\beta_{f,B}(M_B)\circ g^*(b^f(M))
=b^{f'}(g'^*M)\circ\bigl(\beta_{f,A}(M)\bigr)_B,
\tag{9.22}
\]

\[
c^{f'}(g^*K)\circ(e_{f,A}(K))_B
=e_{f,B}(K_B)\circ g'^*c^f(K).
\tag{9.23}
\]

Both identities say that two ways of passing from \((A,Y)\) to \((B,Y')\) give the same map: changing coefficients first and then base, or the other way round. We prove the first by splitting \(Rf_!\) into an open and a proper part, and deduce the second from the first by transposition.

To prove (9.22), choose a compactification \(f=\bar f\circ j\) and pull it back to \(Y'\). Since \(b\) and \(\beta\) for \(f\) are the composites of those for \(j\) and for \(\bar f\), it suffices to treat these two cases. For the open immersion \(j\), extension by zero commutes with pullback and with extension of scalars, and both sides of (9.22) are this identification. For the proper map \(\bar f\), both \(b^{\bar f}\) and proper base change are transposes of the counit of \(L\bar f^*\dashv R\bar f_*\), the first after extension of scalars and the second after pullback; here we work on the derived categories of sheaves of \(A\)-modules on the pro-étale sites. Therefore both sides of (9.22) transpose to the pulled-back counit with scalars extended.

For the second identity, recall that \(c^f\) and \(e_{f,A}\) are defined by transposing the counit composed with \((b^f)^{-1}\), respectively with \(\beta_{f,A}^{-1}\) (see (9.16)). Transposing both sides of (9.23) under \(Rf'_{!,B}\dashv f'^!_B\) thus gives two maps \(Rf'_{!,B}\bigl(g'^*((f_A^!K)_B)\bigr)\to g^*K_B\). Each is built from the counit \(\varepsilon_{f,A}(K)\), the comparisons \(b\) and the base change maps \(\beta\), and the counit commutes with extension of scalars and with pullback. By naturality of \(b\) and \(\beta\), the two maps therefore agree once (9.22) holds.

These two identities let us assemble the finite base change maps into a single map of towers. Its levels are linked by the change-of-coefficients maps along \(\mathcal O_E/\ell^{n+1}\to\mathcal O_E/\ell^n\).

Let \(E/\mathbf Q_\ell\) be a finite extension and work in the category of towers of (9.18). Extension by zero and proper base change define \(\beta_{f,\bullet}:g^*Rf_{!,\bullet}\to Rf'_{!,\bullet}g'^*\), and transposing \(g^*(\varepsilon_{f,\bullet})\circ\beta_{f,\bullet}^{-1}\) under the adjunction on towers for \(f'\) defines \(e_{f,\bullet}:g'^*f_\bullet^!\to f'^!_\bullet g^*\). At level \(n\) these are \(\beta_{f,\mathcal O_E/\ell^n}\) and the exceptional base change map with coefficients \(\mathcal O_E/\ell^n\), so they are isomorphisms; for the second, this is Proposition proposition-exceptional-inverse-image-smooth-morphism with \(\Lambda=\mathcal O_E/\ell^n\). Moreover, the transition maps of \(f_\bullet^!\) are induced by the change-of-coefficients maps \(c\) (Proposition \hyperref[bp58]{9.18}), and by (9.23) the levels of \(e_{f,\bullet}\) commute with them. So \(e_{f,\bullet}\) is an isomorphism of functors on towers. In particular, for a normalized tower \(K_\bullet\) both \(g'^*f_\bullet^!K_\bullet\) and \(f'^!_\bullet g^*K_\bullet\) are normalized.

The functors \(\operatorname{red}\) of Lemma \hyperref[bp55]{9.17} carry the maps over \(\mathcal O_E\) (base change, composition, counits and transposes) to the maps on towers. We check this in steps. First, the functor \(\operatorname{red}\) is a derived pullback of ringed topoi, so it commutes with \(g^*\) by associativity of derived pullbacks, and it commutes with extension by zero because the right adjoints commute. Next, for proper \(\bar f\), the comparison \(\operatorname{red}_YR\bar f_*M\to R\bar f_{*,\bullet}\operatorname{red}_XM\) is the transpose of \(\operatorname{red}_X\) of the counit of \(\bar f^*\dashv R\bar f_*\), and its level \(n\) is the change-of-coefficients map of Lemma \bpout{7.6} for \(\mathcal O_E\to\mathcal O_E/\ell^n\). As evaluations detect isomorphisms, the comparison is an isomorphism for \(M\in\Dbc(X,\mathcal O_E)\). Composing with the compatibility with \(j_!\) we obtain the comparison \(\operatorname{red}_YRf_{!,\mathcal O_E}M\simeq Rf_{!,\bullet}\operatorname{red}_XM\) for a compactification \(f=\bar f\circ j\). The formal argument in the proof of Proposition \bpout{7.10} shows that, under these comparisons, \(\operatorname{red}\) carries the unit and counit of \(\bar f^*\dashv R\bar f_*\) over \(\mathcal O_E\) to the unit and counit on towers. Since the base change and composition isomorphisms of Proposition \bpout{5.28} are composites of such units and counits and of the identifications of pullbacks and of extensions by zero, \(\operatorname{red}\) carries them, with coefficients \(\mathcal O_E\), to \(\beta_{f,\bullet}\) and to the corresponding composition isomorphisms on towers. For the counit, the adjunction of Proposition \hyperref[bp61]{9.19} is transported from the adjunction on towers, so \(\operatorname{red}_Y\) carries \(\varepsilon_f(K)\) to \(\varepsilon_{f,\bullet}(\operatorname{red}_YK)\), once the sources are identified by this comparison and by the isomorphism \(\operatorname{red}_Xf^!_{\mathcal O_E}K\simeq f^!_\bullet\operatorname{red}_YK\) of its proof. Since a transpose is determined by the counit, \(\operatorname{red}\) also carries transposes under \(Rf_{!,\mathcal O_E}\dashv f^!_{\mathcal O_E}\) to transposes on towers. Finally, by Lemma \hyperref[bp55]{9.17}, \(\operatorname{red}\) is an equivalence on \(\Dbc(-,\mathcal O_E)\) whose inverse is \(R\varprojlim\). A morphism of \(\Dbc(-,\mathcal O_E)\) is therefore \(R\varprojlim\) of its image under \(\operatorname{red}\), up to the unit isomorphisms \(M\simeq R\varprojlim\operatorname{red}M\).

Let \(K\in\Dbc(Y,\mathcal O_E)\), and let \(\operatorname{red}_YK=(K/\ell^n)_n\) be its tower of reductions (Lemma \hyperref[bp55]{9.17}). Since inverse image commutes with reduction modulo \(\ell^n\), and \(f_{\mathcal O_E}^!K\simeq R\varprojlim_nf_{\mathcal O_E/\ell^n}^!(K/\ell^n)\) by Proposition \hyperref[bp61]{9.19} and (9.20), both sides below are \(R\varprojlim_n\) of the corresponding terms at level \(n\) for \(K/\ell^n\). Applying \(R\varprojlim\) to \(e_{f,\bullet}\) at \(\operatorname{red}_YK\) we obtain the integral exceptional base change isomorphism

\[
g'^{*}f_{\mathcal O_E}^{!}K
\xrightarrow{\sim}f'^{!}_{\mathcal O_E}g^{*}K;
\tag{9.24}
\]

for \(E=\mathbf Q_\ell\) the coefficients are \(\mathbf Z_\ell\). With coefficients \(E\) and \(\overline{\mathbf Q}_\ell\) we define exceptional base change by inverting \(\ell\) in (9.24) on integral models and then extending scalars from a finite extension of \(\mathbf Q_\ell\). This is possible because \(f^*\), \(f^!\) and \(Rf_!\) with these coefficients are obtained from the integral functors in the same way (Proposition \hyperref[bp66]{9.20} and Theorem \hyperref[bp69]{9.21}).

In each case the map obtained is the exceptional base change map of Definition \hyperref[bp5]{9.1}. For \(\mathcal O_E\), the map of the definition is the transpose under \(Rf'_{!,\mathcal O_E}\dashv f'^!_{\mathcal O_E}\) of \(g^*(\varepsilon_f(K))\circ\beta^{-1}\). By the preceding paragraph, \(\operatorname{red}\) carries this composite to \(g^*(\varepsilon_{f,\bullet})\circ\beta_{f,\bullet}^{-1}\) at \(\operatorname{red}_YK\), and it carries the transpose to the transpose on towers, which is \(e_{f,\bullet}(\operatorname{red}_YK)\). Hence the map of the definition is \(R\varprojlim\) of \(e_{f,\bullet}(\operatorname{red}_YK)\), that is, (9.24). For \(E\) and \(\overline{\mathbf Q}_\ell\) we argue in the same way, with extension of scalars along \(\mathcal O_E\to E\) and \(E\to\overline{\mathbf Q}_\ell\) in place of \(\operatorname{red}\). Indeed, extension of scalars commutes with \(g^*\), extension by zero and \(R\bar f_*\) through the change-of-coefficients maps of Definition \bpout{7.4}, and by Proposition \bpout{7.10} it carries the units and counits of \(\bar f^*\dashv R\bar f_*\) and of \(Rf_!\dashv f^!\) to the corresponding units and counits. So it carries the base change maps for \(Rf_!\), the counits, and therefore the exceptional base change maps over \(\mathcal O_E\) to those over \(E\). In the same way those over a finite extension of \(\mathbf Q_\ell\) go to those over \(\overline{\mathbf Q}_\ell\).

Since the exceptional base change map is an isomorphism for every admissible coefficient ring, we can state that \(\gamma_f\) commutes with it. For \(g\) the inclusion of a closed point \(y\) of \(Y\), this says that \(\gamma_f\) restricts on the fibre \(X_y\) to the isomorphism \(\gamma\) of the smooth \(\kappa(y)\)-scheme \(X_y\).

% object: proposition-compatibility-purity-base-change (blueprint/chap-Smooth-purity/sec-Compatibilities/subsec-Base-change-and-composition/proposition-compatibility-purity-base-change/_section.md)
\phantomsection\label{bp76}%
\begin{bpstatementit}{Proposition~9.22}{ (Compatibility of purity with base change)}
Let \(\Lambda\) be any admissible coefficient ring and let \(K\in D^b_{\mathrm{ctf}}(Y,\Lambda)\). Identify the sources by \(g'^{*}f^{*}\simeq(f')^{*}g^{*}\) and the targets by the exceptional base change map of Definition \hyperref[bp5]{9.1}. Then

\[
g'^{*}\gamma_f(K)
=\gamma_{f'}(g^{*}K)
\quad\text{in }D^b_{\mathrm{ctf}}(X',\Lambda).
\]
\end{bpstatementit}

% proof: proof/proposition-compatibility-purity-base-change
\begin{proof}[{Proof}]\phantomsection\label{bp77}%
The map \(\gamma_f\) has two factors, the orientation of \(\omega_f\) and the map \(\mu_{f,K}\), and we show that each commutes with base change. For the orientation this is a local statement about rank one sheaves, and it follows from the naturality of the deformation to the normal bundle. For \(\mu_{f,K}\) we compare transposes under the adjunction for \(f'\).

For a finite ring \(A\) we compare the orientations and the maps \(\mu_{f,K}\) directly with coefficients in \(A\), and for \(\mathcal O_E\) we compare the maps on towers \(\gamma_{f,\bullet}\) of (9.19) and apply \(R\varprojlim\). The cases \(E\) and \(\overline{\mathbf Q}_\ell\) follow, since \(\gamma_{f,E}\) and \(\gamma_{f,\overline{\mathbf Q}_\ell}\) are obtained from \(\gamma_{f,\mathcal O_E}\) by inverting \(\ell\) and by extension of scalars from a finite extension of \(\mathbf Q_\ell\). We write out the calculation for towers; dropping the index \(\bullet\) gives the calculation at finite level. There the map \(\gamma_{f,\bullet}\) is the composite of the orientation \((\mathcal O_{E,X}/\ell^\bullet)(d)[2d]\xrightarrow{\sim}\omega_{f,\bullet}\) and the map \(f_\bullet^*K_\bullet\otimes\omega_{f,\bullet}\to f_\bullet^!K_\bullet\) for \(K_\bullet=\operatorname{red}_YK\).

\emph{Orientation.} After removing the shift \([2d]\), the two orientations to be compared are maps between sheaves of towers of rank one in degree zero. Since maps between sheaves form a sheaf, and a map of sheaves of towers is a compatible family of maps of sheaves, equality may be tested locally on \(X'\) and at each level \(n\). At a fixed level, the diagonal, its normal bundle, the deformation to the normal bundle and the trace map of vector bundles all commute with base change (Lemma lemma-naturality-under-base-change-etale), and therefore so does the orientation built from them.

\emph{The map \(f^*K\otimes\omega_f\to f^!K\).} This factor is not a map of sheaves, so we cannot compare the two maps locally. However, a map into \(f'^!_\bullet g^*K_\bullet\) is determined by its transpose under the adjunction on towers for \(f'\). By the definition of \(e_{f,\bullet}\), the transpose of the pulled-back map \(\mu_{f,K}\) for \(f\) is the pulled-back counit for \(f\), preceded by the map of the projection formula. Since the projection formula is natural under pullback, the transpose of the map for \(f'\) is the same composite, as in the counit calculation in the proof of Proposition proposition-exceptional-inverse-image-smooth-morphism, carried out with towers.
\end{proof}

% text: prose-09-31
\phantomsection\label{bp78}%
\textbf{Composition.} For composition we again fix the composition isomorphisms for \(f^!\) with every coefficient ring before comparing the maps \(\gamma\). From now on, \(g\) no longer denotes the base change in a square: let \(X\xrightarrow{f}Y\xrightarrow{g}Z\) be smooth of pure relative dimensions \(d\) and \(e\). Composing the finite adjunctions, we identify \((g\circ f)_A^!\) with \(f_A^!g_A^!\), compatibly with \(R(g\circ f)_{!,A}\simeq Rg_{!,A}Rf_{!,A}\), and the counit of the composite adjunction is the composite of the two counits.

The change-of-coefficients maps are compatible with composition: for \(M\in D^b_{\mathrm{ctf}}(X_{\mathrm{pro}\acute{e}t},A)\) and \(K\in D^b_{\mathrm{ctf}}(Z_{\mathrm{pro}\acute{e}t},A)\),

\[
b^{g\circ f}(M)=Rg_{!,B}(b^f(M))\circ b^g(Rf_{!,A}M),
\qquad
c^{g\circ f}(K)=f_B^!c^g(K)\circ c^f(g_A^!K).
\]

The first identity holds on each free term of a resolution of \(B\) over \(A\), since the composition isomorphism is natural, and the totalization in Lemma \hyperref[bp47]{9.15} extends it to arbitrary \(A\to B\). We get the second by transposing under the composite adjunction.

On towers, the composition isomorphism for \(Rf_{!,\bullet}\) is built from the same constructions. The isomorphism \((g\circ f)^!_\bullet\simeq f^!_\bullet g^!_\bullet\) is the isomorphism of right adjoints given by the composite adjunction on towers. Since evaluation at level \(n\) carries the counits on towers to the finite counits (Proposition \hyperref[bp58]{9.18}), the level \(n\) of these isomorphisms is the finite one, and by the identities above they commute with the transition maps, which are built from \(b\) and \(c\). Applying \(R\varprojlim\) gives the integral composition isomorphisms

\[
R(g\circ f)_{!,\mathcal O_E}\simeq Rg_{!,\mathcal O_E}Rf_{!,\mathcal O_E},
\qquad
(g\circ f)_{\mathcal O_E}^{!}\simeq f_{\mathcal O_E}^{!}g_{\mathcal O_E}^{!},
\]

between \(\Dbc(X,\mathcal O_E)\) and \(\Dbc(Z,\mathcal O_E)\); for \(E=\mathbf Q_\ell\) the coefficients are \(\mathbf Z_\ell\). As for base change, inverting \(\ell\) and extending scalars give the composition isomorphisms with coefficients \(E\) and \(\overline{\mathbf Q}_\ell\). These isomorphisms \((g\circ f)^!\simeq f^!g^!\) are the composition isomorphisms of Proposition proposition-composition-open-restriction-compactifications-f. Indeed, the isomorphism of that proposition is characterized by the following property: it carries the counit of \(R(g\circ f)_!\dashv(g\circ f)^!\), composed with the composition isomorphism for \(Rf_!\), to the composite of the counits for \(f\) and \(g\). By uniqueness of adjoints it therefore suffices to check this property for every coefficient ring. With finite coefficients it holds by definition. On towers the composition isomorphism is the one of the composite adjunction, so in \(D(\operatorname{Tow})\) it carries the counit on towers of \(g\circ f\) to the composite of the counits on towers. By the comparison following the construction of \(e_{f,\bullet}\) in Subsection \hyperref[bp74]{9.6.1}, \(\operatorname{red}\) carries the counits and the composition isomorphism for \(Rf_!\) with coefficients \(\mathcal O_E\) to those on towers, and applying \(R\varprojlim\) gives the property over \(\mathcal O_E\). Similarly, the same comparison shows that extension of scalars along \(\mathcal O_E\to E\) and \(E\to\overline{\mathbf Q}_\ell\) carries the counits and the composition isomorphisms for \(Rf_!\) to the corresponding ones, which gives the property over \(E\) and \(\overline{\mathbf Q}_\ell\).

With these composition isomorphisms in place for every coefficient ring, the statement takes the following form. In topological terms it says that the orientation of a composite of submersions is the product of the orientations.

% object: proposition-compatibility-purity-composition (blueprint/chap-Smooth-purity/sec-Compatibilities/subsec-Base-change-and-composition/proposition-compatibility-purity-composition/_section.md)
\phantomsection\label{bp79}%
\begin{bpstatementit}{Proposition~9.23}{ (Compatibility of purity with composition)}
For every admissible \(\Lambda\) and \(K\in D^b_{\mathrm{ctf}}(Z,\Lambda)\), the map \(\gamma_{g\circ f}(K)\) of (9.2) is the composite

\[
\begin{aligned}
(g\circ f)^{*}K(d+e)[2d+2e]
&\simeq f^{*}\bigl(g^{*}K(e)[2e]\bigr)(d)[2d]\\
&\xrightarrow{\ f^{*}(\gamma_g(K))(d)[2d]\ }
f^{*}\bigl(g^!K\bigr)(d)[2d]\\
&\xrightarrow{\ \gamma_f(g^!K)\ }
f^!g^!K
\simeq(g\circ f)^!K
\end{aligned}
\]

in \(D^b_{\mathrm{ctf}}(X,\Lambda)\), where the coefficient \(\Lambda\) is omitted from the displayed composite.
\end{bpstatementit}

% proof: proof/proposition-compatibility-purity-composition
\begin{proof}[{Proof}]\phantomsection\label{bp80}%
As for base change, we compare orientations and the maps \(\mu_{f,K}\), first over any finite admissible ring and then for towers. The orientations are compared on a chart adapted to both \(f\) and \(g\), where each is iterated integration along affine lines, with the coordinates ordered so that the two iterated integrals are the same. Applying \(R\varprojlim\) then gives the integral identity, and inverting \(\ell\) and descending to a finite extension of \(\mathbf Q_\ell\) give the identities for coefficient fields.

\emph{Orientation, finite coefficients.} Both orientations are maps of rank one locally constant sheaves after removing twist and shift, so it suffices to compare them on a cover of \(X\) by charts, and we choose charts adapted to both morphisms. That is, choose an open \(W\subset Y\) with an étale \(Z\)-map \(W\to\mathbf A^e_Z\), and an open \(U\subset f^{-1}(W)\) with an étale \(W\)-map \(U\to\mathbf A^d_W\). Then the composite \(U\to\mathbf A^d_W\to\mathbf A^d_{\mathbf A^e_Z}=\mathbf A^{d+e}_Z\) is étale, and we order the coordinates of \(\mathbf A^{d+e}_Z\) with the \(e\) coordinates of \(g\) first and the \(d\) coordinates of \(f\) last. Over this chart, both orientations are adjoint to iterated integration along affine lines, in the same order: first the \(d\) coordinates of \(f\), then the \(e\) coordinates of \(g\). Since the recursive construction preceding (9.11) integrates the last coordinate first, the integration on \(\mathbf A^{d+e}_Z\) is exactly this composite. Each diagonal orientation agrees on a chart with the integration in its coordinates, so the two orientations agree on the cover.

\emph{The maps \(\mu_{f,K}\), finite coefficients.} The two maps agree because their transposes under the composite adjunction do. Indeed, the transpose of the iterated map applies the projection formula and the counit for \(f\), then the projection formula and the counit for \(g\), whereas the transpose of the map for \(g\circ f\) applies the projection formula for \(g\circ f\) and its counit. The two isomorphisms given by the projection formula agree by associativity of tensor products, and the counit of the composite adjunction is the composite of the two counits.

\emph{Towers.} The same two arguments apply in the category of towers, with its maps \(\gamma_{f,\bullet}\) and its composite adjunction. Since the orientations are maps of sheaves of towers of rank one, we may test their equality at each level, while the maps (9.4) are compared by the same transposition.
\end{proof}


% object: subsec-The-trace-map (blueprint/chap-Smooth-purity/sec-Compatibilities/subsec-The-trace-map/_section.md), content hidden
\subsubsection*{\texorpdfstring{9.6.2\quad The trace map}{9.6.2 The trace map}}
\phantomsection\label{bp81}%
\addcontentsline{toc}{subsubsection}{\texorpdfstring{9.6.2\quad The trace map}{9.6.2 The trace map}}
% text: prose-09-32
\phantomsection\label{bp82}%
On a compact oriented manifold of dimension \(n\), Poincaré duality pairs \(H^i\) with \(H^{n-i}\) by cup product followed by integration \(H^n(M)\to\mathbf R\), which sends the class of a point to \(1\). To obtain the Poincaré pairing of Weil's axiom (c) we need the analogue of integration. By adjunction we already have a map, the counit \(Rf_!f^!K\to K\), but its source involves \(f^!\), which we only know through its universal property. Composing with \(\gamma_f\) replaces \(f^!K\) by the twisted and shifted inverse image, so the result is a map out of compactly supported cohomology of an explicit complex. For example, for \(X\) smooth of dimension \(d\) over an algebraically closed field \(k\) and \(K=\Lambda\), it is a map \(H_c^{2d}(X_{\mathrm{pro}\acute{e}t},\Lambda_X(d))\to\Lambda\). Here the twist cannot be dropped, for the reason explained before Definition \bpout{2.53}: geometric Frobenius acts on \(H^2(\mathbf P^1_k,\mathbf Z_\ell)\) by \(q\), so only the twisted trace map commutes with Frobenius (Section \bpout{10.5}). The trace map is used in the Poincaré pairing (Theorem \bpout{10.12}) and the trace isomorphism in top degree (Theorem \bpout{10.15}).

Let \(f:X\to Y\) be smooth of pure relative dimension \(d\), and let \(\Lambda\) be any admissible coefficient ring. We omit \(\Lambda\) from the notation and let \(\varepsilon_f(K):Rf_!f^!K\to K\) be the counit of the adjunction \(Rf_!\dashv f^!\) between \(D^b_{\mathrm{ctf}}(X,\Lambda)\) and \(D^b_{\mathrm{ctf}}(Y,\Lambda)\).

% object: definition-trace-map (blueprint/chap-Smooth-purity/sec-Compatibilities/subsec-The-trace-map/definition-trace-map/_section.md)
\phantomsection\label{bp83}%
\begin{bpstatement}{Definition~9.24}{ (Trace map)}
For \(K\in D^b_{\mathrm{ctf}}(Y,\Lambda)\), the \emph{trace map} is the composite

\[
\operatorname{Tr}_{f}(K):\quad
Rf_!\bigl(f^{*}K(d)[2d]\bigr)
\xrightarrow{\ Rf_!\gamma_{f}(K)\ }
Rf_!f^!K
\xrightarrow{\ \varepsilon_{f}(K)\ }K.
\tag{9.25}
\]
\end{bpstatement}

% object: remark-trace-map-special-cases (blueprint/chap-Smooth-purity/sec-Compatibilities/subsec-The-trace-map/remark-trace-map-special-cases/_section.md)
\phantomsection\label{bp84}%
\begin{bpstatement}{Remark~9.25}{ (The trace map in two special cases)}
\begin{enumerate}
\item For the projection \(\mathbf A^d_Y\to Y\) and \(K=\Lambda_Y\), the trace map is the trace map of affine space (9.11) with the corresponding coefficients, which on each geometric fibre sends the class of the origin to \(1\).
\item For a finite étale morphism, where \(d=0\), the trace map is summation over the fibres.
\end{enumerate}
\end{bpstatement}

% text: prose-09-33
\phantomsection\label{bp85}%
By the first part of the remark, for affine space the trace map is the explicit integration along lines built from \(\mathbf P^1\). The trace map of a general smooth morphism is glued from such charts, and the next statement says that the gluing respects composition, as integration over the fibres of a composite of submersions is iterated integration.

% object: proposition-transitivity-trace-map (blueprint/chap-Smooth-purity/sec-Compatibilities/subsec-The-trace-map/proposition-transitivity-trace-map/_section.md)
\phantomsection\label{bp86}%
\begin{bpstatementit}{Proposition~9.26}{ (Transitivity of the trace map)}
Let \(X\xrightarrow{f}Y\xrightarrow{g}Z\) be smooth of pure relative dimensions \(d\) and \(e\), and \(K\in D^b_{\mathrm{ctf}}(Z,\Lambda)\) for any admissible \(\Lambda\). Then \(\operatorname{Tr}_{g\circ f}(K)\) is the composite

\[
\begin{gathered}
R(g\circ f)_!
\bigl((g\circ f)^{*}K(d+e)[2d+2e]\bigr)
\simeq Rg_!Rf_!
\bigl(f^{*}
(g^{*}K(e)[2e])(d)[2d]\bigr)\\
\xrightarrow{\ Rg_!
(\operatorname{Tr}_{f}(g^{*}K(e)[2e]))\ }
Rg_!\bigl(g^{*}K(e)[2e]\bigr)
\xrightarrow{\ \operatorname{Tr}_{g}(K)\ }K.
\end{gathered}
\]
\end{bpstatementit}

% proof: proof/proposition-transitivity-trace-map
\begin{proof}[{Proof}]\phantomsection\label{bp87}%
Each trace map is a map \(\gamma\) followed by a counit. The maps \(\gamma\) compose by Proposition \hyperref[bp79]{9.23}, and the counits compose because the counit of a composite adjunction is the composite of the counits, so the only thing to arrange is the order of the factors. If we expand both traces using (9.25), naturality of \(\gamma_f\) lets us move the map \(\gamma_g(K):g^{*}K(e)[2e]\to g^!K\) past the map \(\gamma_f\) of the first trace. The composite then consists of the two maps \(\gamma_f\) and \(\gamma_g\), which together are \(\gamma_{g\circ f}(K)\) by Proposition \hyperref[bp79]{9.23}, followed by

\[
Rg_!Rf_!f^!g^!K
\xrightarrow{\ Rg_!(\varepsilon_{f}(g^!K))\ }
Rg_!g^!K
\xrightarrow{\ \varepsilon_{g}(K)\ }K.
\]

This is the counit of the composite adjunction \(R(g\circ f)_!\dashv(g\circ f)^!\), so the composite is the map \(\gamma_{g\circ f}(K)\) followed by its counit, that is, \(\operatorname{Tr}_{g\circ f}(K)\).
\end{proof}


% object: subsec-Change-of-coefficients (blueprint/chap-Smooth-purity/sec-Compatibilities/subsec-Change-of-coefficients/_section.md), content hidden
\subsubsection*{\texorpdfstring{9.6.3\quad Change of coefficients}{9.6.3 Change of coefficients}}
\phantomsection\label{bp88}%
\addcontentsline{toc}{subsubsection}{\texorpdfstring{9.6.3\quad Change of coefficients}{9.6.3 Change of coefficients}}
% text: prose-09-34
\phantomsection\label{bp89}%
The last compatibility compares the maps \(\gamma_f\) for different coefficient rings. For a map \(A\to B\) of finite rings the comparison \(c_{A,B}\) of Lemma \hyperref[bp47]{9.15} is the change-of-coefficients map for \(f^!\) of Proposition \bpout{7.10}, since both are defined as the transpose of the extended counit composed with the inverse of the comparison for \(Rf_!\).

% object: proposition-compatibility-purity-change-coefficients (blueprint/chap-Smooth-purity/sec-Compatibilities/subsec-Change-of-coefficients/proposition-compatibility-purity-change-coefficients/_section.md)
\phantomsection\label{bp90}%
\begin{bpstatementit}{Proposition~9.27}{ (Compatibility of purity with change of coefficients)}
Let \(f:X\to Y\) be smooth of pure relative dimension \(d\), let \(\Lambda\to\Lambda'\) be a coefficient map of Proposition \bpout{7.10} and let \(K\in D^b_{\mathrm{ctf}}(Y,\Lambda)\). Then the change-of-coefficients isomorphism \((f^!K)_{\Lambda'}\to f^!(K_{\Lambda'})\) carries \((\gamma_{f,\Lambda}(K))_{\Lambda'}\) to \(\gamma_{f,\Lambda'}(K_{\Lambda'})\), the sources being identified by the change-of-coefficients isomorphism for \(f^*\). The same holds for the trace maps: the change-of-coefficients isomorphism for \(Rf_!\) carries \((\operatorname{Tr}_{f,\Lambda}(K))_{\Lambda'}\) to \(\operatorname{Tr}_{f,\Lambda'}(K_{\Lambda'})\).
\end{bpstatementit}

% proof: proof/proposition-compatibility-purity-change-coefficients
\begin{proof}[{Proof}]\phantomsection\label{bp91}%
Both kinds of comparison maps are compatible with composition of coefficient maps, so it suffices to treat the elementary maps in the list of Proposition \bpout{7.10}. For a map \(A\to B\) of finite rings the assertion for \(\gamma_f\) is (9.17), and the assertion for the traces follows from it and the counit identity (9.16). A continuous map \(\mathcal O_E\to A\) factors through \(\mathcal O_E/\ell^n\) for some \(n\), so it suffices to treat \(\mathcal O_E\to\mathcal O_E/\ell^n\). By construction (Subsection \hyperref[bp50]{9.5.2}), \(\operatorname{red}_X\) carries \(\gamma_{f,\mathcal O_E}(K)\) to \(\gamma_{f,\bullet}(\operatorname{red}_YK)\), whose level \(n\) is \(\gamma_{f,\mathcal O_E/\ell^n}(K/\ell^n)\). Here the targets are identified by the isomorphism \(\operatorname{red}_Xf^!_{\mathcal O_E}K\simeq f^!_\bullet\operatorname{red}_YK\), whose level \(n\) is an isomorphism \((f^!_{\mathcal O_E}K)/\ell^n\simeq f^!_{\mathcal O_E/\ell^n}(K/\ell^n)\), and this identification is the change-of-coefficients map of Proposition \bpout{7.10}. Indeed, the integral counit reduces to the finite one, namely to level \(n\) of the counit on towers (Subsection \hyperref[bp74]{9.6.1} and Proposition \hyperref[bp58]{9.18}), and both maps are transposes of that counit composed with the comparison for \(Rf_!\). So this case follows from the finite one. For \(\mathcal O_E\to E\) and \(E\to\overline{\mathbf Q}_\ell\) the assertion holds by construction, because \(\gamma_{f,E}\) and \(\gamma_{f,\overline{\mathbf Q}_\ell}\) were defined by inverting \(\ell\) in the integral adjunction and by extension of scalars from finite extensions of \(\mathbf Q_\ell\) (Subsection \hyperref[bp64]{9.5.3}). For a finite extension \(F/E\) we showed before Theorem \hyperref[bp69]{9.21} that extension of scalars carries the pair of towers for \(\mathcal O_E\) to the pair for \(\mathcal O_F\) by a counit-preserving isomorphism. By applying \(R\varprojlim\) and inverting \(\ell\) we get the cases \(\mathcal O_E\to\mathcal O_F\) and \(E\to F\). In every case the statement for the traces follows from the one for \(\gamma_f\), because the change-of-coefficients maps for \(Rf_!\) and \(f^!\) are compatible with the counits.
\end{proof}



% object: sec-Proof-of-smooth-purity (blueprint/chap-Smooth-purity/sec-Proof-of-smooth-purity/_section.md), content hidden
\subsection*{\texorpdfstring{9.7\quad Proof of smooth purity}{9.7 Proof of smooth purity}}
\phantomsection\label{bp92}%
\addcontentsline{toc}{subsection}{\texorpdfstring{9.7\quad Proof of smooth purity}{9.7 Proof of smooth purity}}
% text: prose-09-35
\phantomsection\label{bp93}%
The proof of Theorem \hyperref[bp9]{9.3} assembles the constructions above; only the trace normalization (1) with integral coefficients needs a separate argument.

% proof: proof/theorem-smooth-purity-all-admissible-coefficients
\begin{proof}[{Proof of Theorem 9.3}]\phantomsection\label{bp94}%
We recall \(\gamma_f\) for each kind of coefficient ring, prove the trace normalization (1) first with finite and then with integral coefficients, and cite (2)–(4).

\emph{The isomorphisms \(\gamma_f\).} For finite \(A\) it is \(\gamma_{f,A}\) of Proposition \hyperref[bp37]{9.12}. For \(\mathcal O_E\) (and \(\mathbf Z_\ell\)) it is \(R\varprojlim\) of the map on towers (9.19), which is an isomorphism with target \(f^!\) by Proposition \hyperref[bp61]{9.19}. For \(E\) (and \(\mathbf Q_\ell\)) and for \(\overline{\mathbf Q}_\ell\) it is the inverse of the counit-preserving isomorphism of Proposition \hyperref[bp41]{9.13} between \(f^!\) and the right adjoints of Proposition \hyperref[bp66]{9.20} and Theorem \hyperref[bp69]{9.21}. In every case, therefore, \(\gamma_f\) is natural in \(K\).

\emph{Property (1), finite coefficients.} Here (1) holds because the orientation obtained from the diagonal agrees on the chart \(\mathbf A^1_Y\) with the one given by the adjunction in coordinates, whose counit at the constant sheaf is \(\operatorname{Tr}_p\) (the comparison following Proposition \hyperref[bp34]{9.11}).

\emph{Property (1), integral and rational coefficients.} For \(\Lambda=\mathcal O_E\) we write the left side \(\phi\) of (1) as \(u\circ\operatorname{Tr}_p\) and show that \(u=1\) by reducing modulo \(\ell^n\).

First, the trace map \(\operatorname{Tr}_p\) of (9.1) is an isomorphism, and this is what makes \(u\) well defined. Its reduction modulo \(\ell\) is the trace map with coefficients \(\mathcal O_E/\ell\), because the trace of Proposition \bpout{6.22} lifts the finite ones, and this reduction is an isomorphism. Indeed, it is obtained by extension of scalars along the free map \(\mathbf Z/\ell\to\mathcal O_E/\ell\) from the trace map with coefficients \(\mathbf Z/\ell\), which is an isomorphism (Section \hyperref[bp15]{9.3}). So the cone \(C\) of \(\operatorname{Tr}_p\) has \(C/\ell=0\), and hence \(C/\pi_E=(C/\ell)\otimes^{\mathbf L}_{\mathcal O_E/\ell}\mathcal O_E/\pi_E=0\). Since the source and target of \(\operatorname{Tr}_p\) lie in \(\Dbc(Y,\mathcal O_E)\), the complex \(C\) is derived complete, and derived Nakayama (Lemma \bpout{2.62}) gives \(C=0\). Thus \(\phi=u\circ\operatorname{Tr}_p\) for a unique endomorphism \(u\) of \(\mathcal O_{E,Y}\), that is, for an element \(u\in H^0(Y,\mathcal O_E)=\varprojlim_nH^0(Y,\mathcal O_E/\ell^n)\).

Next, we show that \(u\equiv1\) modulo \(\ell^n\) for every \(n\), so that \(u=1\). The trace part of Proposition \hyperref[bp90]{9.27} for \(\mathcal O_E\to\mathcal O_E/\ell^n\) shows that the reduction of \(\phi\) modulo \(\ell^n\) is the left side of (1) with coefficients \(\mathcal O_E/\ell^n\), which by the finite case equals the trace map with these coefficients. This is also the reduction of \(\operatorname{Tr}_p\), and that reduction is an isomorphism. Hence \(u\equiv1\) modulo \(\ell^n\). For \(E\) and \(\overline{\mathbf Q}_\ell\) both sides are obtained from the integral ones by inverting \(\ell\) and extending scalars.

\emph{Properties (2)–(4).} These are Propositions \hyperref[bp76]{9.22}, \hyperref[bp79]{Proposition 9.23} and \hyperref[bp90]{Proposition 9.27}.
\end{proof}

% text: prose-09-36
\phantomsection\label{bp95}%
For the structure morphism of a smooth scheme, Theorem \hyperref[bp9]{9.3} computes the dualizing complex used in Poincaré duality (Chapter \bpout{10}).

% object: corollary-dualizing-complex-smooth-scheme (blueprint/chap-Smooth-purity/sec-Proof-of-smooth-purity/corollary-dualizing-complex-smooth-scheme/_section.md)
\phantomsection\label{bp96}%
\begin{bpstatementit}{Corollary~9.28}{ (The dualizing complex of a smooth scheme)}
Let \(a:X\to\operatorname{Spec}k\) be smooth of pure dimension \(d\), and let \(\Lambda\) be any admissible coefficient ring. There is a natural isomorphism \(a^!\Lambda_{\operatorname{Spec}k}\simeq\Lambda_X(d)[2d]\) in \(D^b_{\mathrm{ctf}}(X,\Lambda)\).
\end{bpstatementit}

% proof: proof/corollary-dualizing-complex-smooth-scheme
\begin{proof}[{Proof}]\phantomsection\label{bp97}%
Apply Theorem \hyperref[bp9]{9.3} with coefficients \(\Lambda\) to \(K=\Lambda_{\operatorname{Spec}k}\), and use \(a^*\Lambda_{\operatorname{Spec}k}=\Lambda_X\).
\end{proof}

% text: prose-09-37
\phantomsection\label{bp98}%
As a check, take \(X=\mathbf A^1_k\) with \(k\) algebraically closed. The adjunction gives \(\Hom(R\Gamma_c(\mathbf A^1,\Lambda),\Lambda[-2])\simeq\Hom(\Lambda_X,a^!\Lambda[-2])
=H^0(\mathbf A^1,\Lambda(1))=\Lambda(1)\), and indeed \(R\Gamma_c(\mathbf A^1,\Lambda)=\Lambda(-1)[-2]\).

If \(k\) is algebraically closed, taking global sections turns the trace \(\operatorname{Tr}_a(\Lambda_{\operatorname{Spec}k})\) into a map \(R\Gamma_c(X_{\mathrm{pro}\acute{e}t},\Lambda_X(d))[2d]\to\Lambda\), and taking \(H^0\) of it gives the trace map \(H_c^{2d}(X_{\mathrm{pro}\acute{e}t},\Lambda_X(d))\to\Lambda\). For example, for \(X=\mathbf P^1_k\) it sends \(c_1(\mathcal O(1))\) to \(1\).



\clearpage\phantomsection\hypertarget{bp-outside}{}%
\section*{Outside this export}
\addcontentsline{toc}{section}{Outside this export}
References marked \ensuremath{{}^\ast} are to parts of the blueprint this document leaves out:
\begin{description}
\item[Corollary 2.29] Continuity of cohomology, in Ⅰ The six operations (\bpslug{corollary-{}\allowbreak{}continuity-{}\allowbreak{}cohomology})
\item[Theorem 2.38] Comparison of direct images, in Ⅰ The six operations (\bpslug{theorem-{}\allowbreak{}comparison-{}\allowbreak{}direct-{}\allowbreak{}images})
\item[Corollary 2.44] Finite Tor amplitude under comparison, in Ⅰ The six operations (\bpslug{corollary-{}\allowbreak{}finite-{}\allowbreak{}tor-{}\allowbreak{}amplitude-{}\allowbreak{}under-{}\allowbreak{}comparison})
\item[Definition 2.53] Tate twists, in Ⅰ The six operations (\bpslug{definition-{}\allowbreak{}tate-{}\allowbreak{}twists})
\item[Lemma 2.62] Reduction and derived Nakayama, in Ⅰ The six operations (\bpslug{lemma-{}\allowbreak{}reduction-{}\allowbreak{}derived-{}\allowbreak{}nakayama})
\item[Definition 3.1] Constructible sheaves, in Ⅰ The six operations (\bpslug{definition-{}\allowbreak{}constructible-{}\allowbreak{}sheaves})
\item[Proposition 3.11] The integral constructibility criterion, in Ⅰ The six operations (\bpslug{proposition-{}\allowbreak{}integral-{}\allowbreak{}constructibility-{}\allowbreak{}criterion})
\item[Lemma 3.13] Flat models and generators, in Ⅰ The six operations (\bpslug{lemma-{}\allowbreak{}flat-{}\allowbreak{}models-{}\allowbreak{}generators})
\item[Proposition 3.27] Inverting $\ell$, in Ⅰ The six operations (\bpslug{proposition-{}\allowbreak{}inverting-{}\allowbreak{}ell})
\item[Theorem 5.2] The exceptional inverse image, in Ⅰ The six operations (\bpslug{theorem-{}\allowbreak{}exceptional-{}\allowbreak{}inverse-{}\allowbreak{}image})
\item[Proposition 5.15] Comparison with the étale derived category, in Ⅰ The six operations (\bpslug{proposition-{}\allowbreak{}comparison-{}\allowbreak{}etale-{}\allowbreak{}derived-{}\allowbreak{}category})
\item[Lemma 5.19] Structure of integral constructible complexes, in Ⅰ The six operations (\bpslug{lemma-{}\allowbreak{}structure-{}\allowbreak{}integral-{}\allowbreak{}constructible-{}\allowbreak{}complexes})
\item[Proposition 5.28] Composition and base change for $Rf_!$, in Ⅰ The six operations (\bpslug{proposition-{}\allowbreak{}composition-{}\allowbreak{}base-{}\allowbreak{}change-{}\allowbreak{}rf-{}\allowbreak{}shriek})
\item[Theorem 5.31] Existence of the exceptional inverse image, in Ⅰ The six operations (\bpslug{theorem-{}\allowbreak{}existence-{}\allowbreak{}exceptional-{}\allowbreak{}inverse-{}\allowbreak{}image})
\item[Section 6.1] Conventions, in Ⅰ The six operations (\bpslug{sec-{}\allowbreak{}Conventions-{}\allowbreak{}2})
\item[Lemma 6.3] The cycle class of a Cartier divisor, in Ⅰ The six operations (\bpslug{lemma-{}\allowbreak{}cycle-{}\allowbreak{}class-{}\allowbreak{}cartier-{}\allowbreak{}divisor})
\item[Lemma 6.4] Cohomology of the projective line and the trace, in Ⅰ The six operations (\bpslug{lemma-{}\allowbreak{}cohomology-{}\allowbreak{}projective-{}\allowbreak{}line-{}\allowbreak{}trace})
\item[Proposition 6.7] The adjunction for the projective line, in Ⅰ The six operations (\bpslug{proposition-{}\allowbreak{}adjunction-{}\allowbreak{}projective-{}\allowbreak{}line})
\item[Proposition 6.9] The exceptional inverse image of a smooth morphism, in Ⅰ The six operations (\bpslug{proposition-{}\allowbreak{}exceptional-{}\allowbreak{}inverse-{}\allowbreak{}image-{}\allowbreak{}smooth-{}\allowbreak{}morphism})
\item[Theorem 6.20] $f^!$ with finite coefficients, in Ⅰ The six operations (\bpslug{theorem-{}\allowbreak{}f-{}\allowbreak{}shriek-{}\allowbreak{}finite-{}\allowbreak{}coefficients})
\item[Proposition 6.22] The $\ell$-adic cycle class of the diagonal, in Ⅰ The six operations (\bpslug{proposition-{}\allowbreak{}ell-{}\allowbreak{}adic-{}\allowbreak{}cycle-{}\allowbreak{}class-{}\allowbreak{}diagonal})
\item[Proposition 7.1] The six operations on $\Dctf$, in Ⅰ The six operations (\bpslug{proposition-{}\allowbreak{}six-{}\allowbreak{}operations-{}\allowbreak{}dctf})
\item[Section 7.2] Change of coefficients, in Ⅰ The six operations (\bpslug{sec-{}\allowbreak{}Change-{}\allowbreak{}of-{}\allowbreak{}coefficients})
\item[Definition 7.4] Change-of-coefficients maps, in Ⅰ The six operations (\bpslug{definition-{}\allowbreak{}change-{}\allowbreak{}of-{}\allowbreak{}coefficients-{}\allowbreak{}maps})
\item[Lemma 7.6] Change of coefficients for direct image, in Ⅰ The six operations (\bpslug{lemma-{}\allowbreak{}change-{}\allowbreak{}coefficients-{}\allowbreak{}direct-{}\allowbreak{}image})
\item[Proposition 7.10] The six operations and change of coefficients, in Ⅰ The six operations (\bpslug{proposition-{}\allowbreak{}six-{}\allowbreak{}operations-{}\allowbreak{}change-{}\allowbreak{}coefficients})
\item[Proposition 7.11] $\Dctf(X,\Qbar)$ as a $2$-colimit, in Ⅰ The six operations (\bpslug{proposition-{}\allowbreak{}dctf-{}\allowbreak{}x-{}\allowbreak{}qbar-{}\allowbreak{}2-{}\allowbreak{}colimit})
\item[Section 7.3] Reduction to finite coefficients, in Ⅰ The six operations (\bpslug{sec-{}\allowbreak{}Reduction-{}\allowbreak{}to-{}\allowbreak{}finite-{}\allowbreak{}coefficients})
\item[Chapter 10] Verdier and Poincaré duality, in Ⅱ Duality and the trace formula (\bpslug{chap-{}\allowbreak{}Verdier-{}\allowbreak{}and-{}\allowbreak{}Poincare-{}\allowbreak{}duality})
\item[Theorem 10.12] Poincaré duality, in Ⅱ Duality and the trace formula (\bpslug{theorem-{}\allowbreak{}poincare-{}\allowbreak{}duality})
\item[Theorem 10.15] The trace isomorphism in top degree, in Ⅱ Duality and the trace formula (\bpslug{theorem-{}\allowbreak{}trace-{}\allowbreak{}isomorphism-{}\allowbreak{}top-{}\allowbreak{}degree})
\item[Section 10.5] Weil complexes, in Ⅱ Duality and the trace formula (\bpslug{sec-{}\allowbreak{}Weil-{}\allowbreak{}complexes})
\item[Proposition 10.20] Purity and trace commute with Frobenius, in Ⅱ Duality and the trace formula (\bpslug{proposition-{}\allowbreak{}purity-{}\allowbreak{}trace-{}\allowbreak{}commute-{}\allowbreak{}frobenius})
\end{description}
\clearpage\phantomsection%
\section*{References}
\addcontentsline{toc}{section}{References}
\bpbibitem{BBD82}{bp100}{A. Beilinson, J. Bernstein, and P. Deligne. \emph{Faisceaux pervers}, volume 100 of \emph{Astérisque}. 1982.}

\bpbibitem{BS15}{bp101}{B. Bhatt and P. Scholze. The pro-étale topology for schemes. \emph{Astérisque}, 369:99–201, 2015. \href{https://people.mpim-bonn.mpg.de/scholze/proetale.pdf}{Author's manuscript}.}

\bpbibitem{Con07}{bp102}{B. Conrad. Deligne's notes on Nagata compactifications. \emph{Journal of the Ramanujan Mathematical Society}, 22:205–257, 2007. \href{https://math.stanford.edu/\~conrad/papers/nagatafinal.pdf}{Author's manuscript}.}

\bpbibitem{Del74}{bp103}{P. Deligne. La conjecture de Weil. I. \emph{Publ. Math. IHÉS}, 43:273–307, 1974.}

\bpbibitem{Del80}{bp104}{P. Deligne. La conjecture de Weil. II. \emph{Publ. Math. IHÉS}, 52:137–252, 1980.}

\bpbibitem{Den22}{bp105}{C. Deninger. There is no “Weil-”cohomology theory with real coefficients for arithmetic curves. \emph{Ann. Sc. Norm. Super. Pisa Cl. Sci.}, 2022. To appear. arXiv:\href{https://arxiv.org/abs/2204.02714}{2204.02714}.}

\bpbibitem{Deu41}{bp106}{M. Deuring. Die Typen der Multiplikatorenringe elliptischer Funktionenkörper. \emph{Abh. Math. Sem. Univ. Hamburg}, 14:197–272, 1941.}

\bpbibitem{Eke90}{bp107}{T. Ekedahl. On the adic formalism. In \emph{The Grothendieck Festschrift, Vol. II}, volume 87 of \emph{Progress in Mathematics}, pages 197–218. Birkhäuser, 1990.}

\bpbibitem{Has36}{bp108}{H. Hasse. Zur Theorie der abstrakten elliptischen Funktionenkörper. I, II, III. \emph{J. Reine Angew. Math.}, 175:55–62, 69–88, 193–208, 1936.}

\bpbibitem{Hum72}{bp109}{J. E. Humphreys. \emph{Introduction to Lie Algebras and Representation Theory}, volume 9 of \emph{Graduate Texts in Mathematics}. Springer, 1972.}

\bpbibitem{Ill94}{bp110}{L. Illusie. Autour du théorème de monodromie locale. \emph{Astérisque}, 223:9–57, 1994. §4. \href{https://www.numdam.org/item/AST_1994__223__9_0/}{Original article}.}

\bpbibitem{Kat88}{bp111}{N. M. Katz. Travaux de Laumon. \emph{Astérisque}, 161–162:105–132, 1988. Séminaire Bourbaki, exp. 691.}

\bpbibitem{KW01}{bp112}{R. Kiehl and R. Weissauer. \emph{Weil conjectures, perverse sheaves and \(\ell\)-adic Fourier transform}, volume 42 of \emph{Ergebnisse der Mathematik}. Springer, 2001.}

\bpbibitem{Lau87}{bp113}{G. Laumon. Transformation de Fourier, constantes d'équations fonctionnelles et conjecture de Weil. \emph{Publ. Math. IHÉS}, 65:131–210, 1987.}

\bpbibitem{Mil86a}{bp114}{J. S. Milne. \href{https://www.jmilne.org/math/xnotes/AVs.pdf}{Abelian Varieties}. In \emph{Arithmetic Geometry}. 1986. Corrected online version.}

\bpbibitem{Mil86b}{bp115}{J. S. Milne. \href{https://www.jmilne.org/math/xnotes/JVs.pdf}{Jacobian Varieties}. In \emph{Arithmetic Geometry}. 1986. Corrected online version.}

\bpbibitem{Ran39}{bp116}{R. A. Rankin. Contributions to the theory of Ramanujan's function \(\tau(n)\) and similar arithmetical functions. II. \emph{Proc. Cambridge Philos. Soc.}, 35:357–372, 1939.}

\bpbibitem{SGA4}{bp99}{M. Artin, A. Grothendieck, and J.-L. Verdier, editors. \emph{Théorie des topos et cohomologie étale des schémas}, volumes 269, 270, 305 of \emph{Lecture Notes in Mathematics}. Springer, 1972–1973. SGA 4.}

\bpbibitem{SGA4½}{bp117}{P. Deligne. \emph{Cohomologie étale}, volume 569 of \emph{Lecture Notes in Mathematics}. Springer, 1977. SGA 4½; \emph{Théorèmes de finitude}, §3, and \emph{Rapport sur la formule des traces}, §3, and \emph{Sommes trigonométriques}, §1.}

\bpbibitem{SGA5}{bp118}{A. Grothendieck, editor. \emph{Cohomologie \(\ell\)-adique et fonctions \(L\)}, volume 589 of \emph{Lecture Notes in Mathematics}. Springer, 1977. SGA 5.}

\bpbibitem{SGA7II}{bp119}{P. Deligne and N. Katz, editors. \emph{Groupes de monodromie en géométrie algébrique}, volume 340 of \emph{Lecture Notes in Mathematics}. Springer, 1973. SGA 7 II; Exposé XIII, (2.1.7).}

\bpbibitem{Stacks}{bp120}{The Stacks Project Authors. \href{https://stacks.math.columbia.edu}{The Stacks Project}. Individual results are cited by their stable tags in the text.}

\bpbibitem{Var05}{bp121}{Y. Varshavsky. Lefschetz–Verdier trace formula and a generalization of a theorem of Fujiwara. 2005. Version 2. arXiv:\href{https://arxiv.org/abs/math/0505564}{math/0505564}. \href{https://arxiv.org/abs/math/0505564v2}{Author's preprint}.}

\bpbibitem{Wei48a}{bp122}{A. Weil. \emph{Sur les courbes algébriques et les variétés qui s'en déduisent}. Hermann, Paris, 1948.}

\bpbibitem{Wei48b}{bp123}{A. Weil. \emph{Variétés abéliennes et courbes algébriques}. Hermann, Paris, 1948.}

\bpbibitem{Wei49}{bp124}{A. Weil. Numbers of solutions of equations in finite fields. \emph{Bull. Amer. Math. Soc.}, 55:497–508, 1949.}

\end{document}
